
\documentclass[11pt,a4paper]{article}
\usepackage[a4paper,margin=2cm]{geometry}
\usepackage{fontspec}
\usepackage{xepersian}
\usepackage{hyperref}
\usepackage{enumitem}
\usepackage{longtable}
\usepackage{array}
\usepackage{xcolor}
\usepackage{titlesec}
\usepackage{fancyhdr}
\usepackage{booktabs}
\usepackage{tabularx}
\usepackage{framed}
\settextfont{DejaVu Sans}
\setlatintextfont{DejaVu Sans}
\setmonofont{DejaVu Sans Mono}
\hypersetup{colorlinks=true, linkcolor=blue!55!black, urlcolor=blue!55!black}
\definecolor{Accent}{HTML}{1F4FD8}
\definecolor{SoftGray}{HTML}{F4F6FA}
\definecolor{DarkGray}{HTML}{333333}
\titleformat{\section}{\Large\bfseries\color{Accent}}{\thesection}{1em}{}
\titleformat{\subsection}{\large\bfseries\color{DarkGray}}{\thesubsection}{1em}{}
\titleformat{\subsubsection}{\normalsize\bfseries\color{DarkGray}}{\thesubsubsection}{1em}{}
\setlength{\headheight}{15pt}
\pagestyle{fancy}
\fancyhf{}
\rhead{گزارش کامل بررسی و توسعه TAFlow}
\lhead{TAFlow Review}
\cfoot{\thepage}
\setlist[itemize]{topsep=3pt,itemsep=2pt,parsep=0pt}
\setlist[enumerate]{topsep=3pt,itemsep=2pt,parsep=0pt}
\renewcommand{\arraystretch}{1.35}
\begin{document}
\begin{titlepage}
\centering
\vspace*{2.5cm}
{\Huge\bfseries گزارش کامل بررسی، اصلاحات، تست و نقشه توسعه TAFlow\par}
\vspace{0.6cm}
{\Large شامل چک لیست اصلاحی، ۳۰۰ قابلیت پیشنهادی، مراحل پیاده سازی و پرامپت جامع Claude\par}
\vspace{1.5cm}
\begin{framed}
\centering
\Large پروژه: \lr{Ramin-Buzarpur/TAFlow}\par
\vspace{0.2cm}
\normalsize دامنه بررسی: معماری، دیتابیس، امنیت، RBAC، workflow، Gradebook، تکلیف، نظرسنجی، فایل، گزارش، UI/UX، تست و استقرار
\end{framed}
\vfill
{\large تاریخ تهیه: ۲۰۲۶-۰۷-۰۳\par}
\vspace{0.4cm}
{\small این سند برای استفاده مستقیم به عنوان دستور کار توسعه و همچنین پرامپت قابل ارسال به Claude آماده شده است.}
\end{titlepage}
\tableofcontents
\newpage
\section{دامنه و وضعیت فعلی پروژه}
این بخش خلاصه دقیق وضعیت فعلی پروژه و محدوده بررسی است.
\begin{itemize}
\item ریپازیتوری Ramin-Buzarpur/TAFlow یک پروژه فول استک برای مدیریت TA و Head TA است و در سطح MVP سنگین قرار دارد، نه فقط یک قالب ظاهری.
\item ساختار اصلی پروژه شامل src، prisma، tests، docs، docker-compose.yml، pnpm-lock.yaml، playwright.config.ts و فایل های مستندات/Prompt است.
\item براساس README، تکنولوژی های اصلی شامل Next.js App Router، TypeScript، PostgreSQL، Prisma، Auth.js/NextAuth، Zod، Argon2id، TOTP، Redis-backed rate limiting، AuditLog، RTL UI، Vitest، Playwright و Docker Compose است.
\item در src/app/api مسیرهای گسترده ای برای account، admin، announcements، auth، calendar، certificates، course-offerings، dashboard، evaluations، exports، files، gradebook، messages، notifications، polls، regrade-requests، reports، search، sessions، surveys، ta-applications، ta-opportunities، talent-pool و tasks وجود دارد.
\item در src/app صفحه های فرانت برای admin، announcements، applications، certificates، courses، dashboard، evaluations، files، gradebook، grades، login، messages، opportunities، register، reports، sessions، settings، surveys و talent-pool وجود دارد.
\item پروژه unit test و e2e test دارد، اما با توجه به تعداد نقش ها و مسیرهای حساس، باید پوشش تست امنیتی و workflow عمیق تر شود.
\item این گزارش یک review ایستا از ساختار عمومی ریپو و محتوای منتشر شده است؛ اجرای کامل pnpm install/build/test در این محیط انجام نشده، پس تمام موارد اجرایی باید در سیستم توسعه پروژه validate شوند.
\end{itemize}

\subsection{منابع بررسی}
\begin{itemize}
\item مخزن اصلی GitHub: \url{https://github.com/Ramin-Buzarpur/TAFlow}
\item README پروژه: \url{https://github.com/Ramin-Buzarpur/TAFlow/blob/main/README.md}
\item مسیرهای API و Frontend در شاخه \lr{src/app} و \lr{src/app/api}
\item schema و migrationها در شاخه \lr{prisma}
\item تست ها در شاخه \lr{tests/unit} و \lr{tests/e2e}
\end{itemize}
\subsection{جمع بندی فنی اولیه}
\begin{itemize}
\item پروژه از نظر ساختار کلی خوب و گسترده است؛ APIها، صفحات، تست ها، migrationها و طراحی RTL وجود دارند.
\item اما پروژه هنوز باید با نگاه QA، Product، Security و Production Hardening بررسی و تکمیل شود.
\item حساس ترین بخش ها عبارت اند از: RBAC، CourseRoleAssignment، TA application workflow، Gradebook permissions، Survey/Poll anonymity، File access security، Certificate verification/revocation، E2E flows for each role و Build/typecheck/deployment.
\end{itemize}
\section{چک لیست کامل اصلاحات، تغییرات و تست های ضروری}
در این بخش ۱۰۰ مورد اصلاحی و تستی اصلی آمده است. این موارد باید قبل از تحویل نهایی پروژه یا قبل از اضافه کردن قابلیت های بزرگ جدید انجام شوند.
\subsection{Package / نصب / سازگاری نسخه ها}
\begin{enumerate}[label=\arabic*.]
\item نسخه های اصلی و حساس باید کامل pin شوند. هر dependency که با \textasciicircum{} یا range باز تعریف شده، مخصوصا Next, React, Prisma, next-auth, argon2, motion, nodemailer, AWS SDK و tooling، باید بررسی و در صورت نیاز ثابت شود.
\item Prisma باید روی 6.19.3 قفل بماند یا پروژه رسما و کامل به Prisma 7 migrate شود. حالت نیمه کاره باعث خطای datasource property url is no longer supported می شود.
\item مسئله pnpm approve-builds باید در README، setup و troubleshooting واضح شود؛ مخصوصا برای @prisma/client, prisma engines, argon2, sharp و esbuild.
\item اسکریپت های آماده مثل pnpm setup:windows, pnpm setup:reset و pnpm setup:prod-check اضافه شود تا اجرای پروژه برای ارائه و توسعه راحت تر شود.
\item فقط pnpm-lock.yaml باید منبع lock رسمی باشد. وجود package-lock.json یا lockfile موازی باید ممنوع شود.
\end{enumerate}
\subsection{Prisma / دیتابیس / migration}
\begin{enumerate}[label=\arabic*. ,resume]
\item تفکیک Course از CourseOffering حفظ شود و در تمام سرویس ها رعایت شود؛ Course تعریف ثابت درس است و CourseOffering ارائه همان درس در یک ترم خاص.
\item CourseRoleAssignment باید تاریخچه assign/revoke/reassign را نگه دارد. partial unique index فعال بودن نقش باید در migration واقعی PostgreSQL تست شود.
\item همه migrationهای init، integrity\_constraints، product\_workflows و integrity\_gaps\_2 باید از صفر روی دیتابیس خام اجرا و validate شوند.
\item برای تمام raw SQL constraintها integration test نوشته شود؛ comment داخل schema کافی نیست.
\item برای GradeRecord.score باید DB-level check وجود داشته باشد: score >= 0 و score <= maxScore؛ مگر اینکه bonus policy جدا تعریف شود.
\item برای Interview و OfficeHourSession باید constraint یا validation قطعی وجود داشته باشد که startsAt < endsAt باشد.
\item برای Semester باید startsAt < endsAt enforce شود.
\item برای TAOpportunity.deadline باید validate شود که deadline بعد از publish و قبل از close منطقی باشد.
\item برای PollVote در حالت anonymous باید respondentHash و partial unique index با race-condition test تست شود.
\item برای SurveyAnswer باید تضمین شود هر respondent فقط یک پاسخ برای هر question ثبت کند.
\item برای ProfessorEvaluation و TAEvaluation باید constraint محدوده امتیاز، مثلا 1 تا 5 یا 1 تا 10، مشخص و enforce شود.
\item برای CertificateRequest باید duplicate request فعال ممنوع باشد، ولی request جدید بعد از rejected/revoked طبق سیاست سیستم امکان پذیر باشد.
\item برای CourseEnrollment باید drop/re-enroll history واقعا کار کند و active enrollment با enrollment تاریخچه ای اشتباه گرفته نشود.
\item برای UploadedFile باید soft-delete policy دقیق شود: چه زمانی فایل فقط hidden می شود و چه زمانی واقعا از storage پاک می شود.
\item seed باید idempotent شود؛ اجرای چندباره pnpm db:seed نباید duplicate data یا crash ایجاد کند.
\end{enumerate}
\subsection{Auth / Session / Security}
\begin{enumerate}[label=\arabic*. ,resume]
\item login، signup، email verification، forgot/reset password، Argon2id، lockout، rate limit، TOTP و groundwork برای Keycloak باید e2e test واقعی داشته باشند.
\item کاربر با وضعیت PENDING\_EMAIL نباید به بخش های اصلی سیستم دسترسی داشته باشد.
\item کاربر SUSPENDED یا DELETED نباید بتواند با session قدیمی وارد پنل شود.
\item بعد از تغییر password باید sessionهای قبلی invalidate شوند.
\item reset password token باید single-use، زمان دار و audit شده باشد.
\item TOTP recovery codes باید اضافه شود یا صریحا تصمیم گرفته شود که وجود ندارند؛ برای production بهتر است وجود داشته باشند.
\item برای Professor، Education Admin و Head TA بهتر است 2FA اجباری شود.
\item rate limit فقط روی login کافی نیست؛ باید register، forgot-password، survey-answer، poll-vote، upload، message، export و certificate هم محدود شوند.
\item در production باید Redis rate limit اجباری باشد؛ fallback in-memory فقط برای dev قابل قبول است.
\item SecurityEvent باید login failed/success، password reset، 2FA fail، permission denied و suspicious export را کامل ثبت کند.
\item Security headers اضافه شود: CSP، HSTS، X-Frame-Options، X-Content-Type-Options، Referrer-Policy و Permissions-Policy.
\item همه mutation APIها شامل POST/PATCH/DELETE باید origin check و CSRF defense مناسب داشته باشند.
\item Keycloak/SSO باید domain دانشگاهی را enforce کند.
\item اگر SSO فعال شد، mapping بین SSO profile و User/Profile باید دقیق و تست شده باشد.
\item audit log باید برای همه عملیات حساس اجباری باشد و bypass مستقیم Prisma بدون audit ممنوع شود.
\end{enumerate}
\subsection{RBAC / Permission}
\begin{enumerate}[label=\arabic*. ,resume]
\item هر permission باید از server و active CourseRoleAssignment خوانده شود، نه از client state.
\item Student نباید بتواند application خودش را approve یا status آن را دستکاری کند.
\item TA نباید grade item غیر assigned را edit کند.
\item Head TA نباید بتواند role Professor یا Admin بدهد.
\item Professor فقط در courseOffering خودش باید دسترسی کامل داشته باشد، نه کل سیستم.
\item Education Admin باید سطح university/department داشته باشد و تمام کارهای حساس او audit شود.
\item revoked role نباید هیچ دسترسی مبتنی بر course داشته باشد.
\item expired role نباید دسترسی داشته باشد.
\item deleted user نباید در queryهای active role دیده شود.
\item command palette باید permission-aware بماند و هیچ نتیجه غیرمجاز نشان ندهد.
\item همه APIهای admin باید role guard صریح و server-side داشته باشند.
\item همه export APIها باید permission check و audit داشته باشند.
\item file download باید permission check و audit داشته باشد.
\item evaluation results باید minimum response count و privacy rule داشته باشد.
\item survey anonymous mode نباید هویت پاسخ دهنده را از طریق UI، API یا export لو بدهد.
\end{enumerate}
\subsection{TA / Head TA Workflow}
\begin{enumerate}[label=\arabic*. ,resume]
\item برای وضعیت های SUBMITTED, UNDER\_REVIEW, SHORTLISTED, INTERVIEW\_INVITED, ACCEPTED, REJECTED, WITHDRAWN باید transition matrix دقیق تعریف شود.
\item نباید از REJECTED مستقیم به ACCEPTED رفت، مگر با admin override audit شده.
\item بعد از deadline نباید application جدید ثبت شود.
\item اگر opportunity closed یا cancelled شد، application جدید باید ممنوع شود.
\item اگر application accepted شد، CourseRoleAssignment باید به صورت transactionally ساخته شود.
\item اگر application برای Head TA پذیرفته شد، role HEAD\_TA باید ساخته شود.
\item سیاست Head TA باید مشخص شود: فقط review/recommendation یا approve/reject نهایی.
\item capacity و requiredTAs باید در پذیرش نهایی رعایت شود.
\item تداخل زمانی با کلاس، امتحان و جلسه باید در انتخاب TA بررسی شود.
\item Professor باید بتواند shortlist بسازد و تاریخچه shortlist ثبت شود.
\item interview schedule، score، notes و result باید کامل ثبت شود.
\item rejection reason دانشجو پسند و internal note باید جدا باشند.
\item appeal/review request برای رد شدن اضافه شود.
\item recommendation engine برای پیشنهاد بهترین TAها اضافه شود.
\item talent pool باید براساس سابقه، نمره، مهارت و availability کامل شود.
\end{enumerate}
\subsection{Gradebook / Assignment / Homework}
\begin{enumerate}[label=\arabic*. ,resume]
\item بخش gradebook شامل weighted categories، grade items، edit history، publish، student private view، CSV export و bulk XLSX import باید تست عمیق شود.
\item مجموع weightها نباید از 100 درصد بیشتر شود و بهتر است سیاست 100 دقیق یا کمتر از 100 مشخص شود.
\item import نمرات باید preview داشته باشد و قبل از نوشتن در DB تمام خطاها را نشان دهد.
\item import باید rollback کامل داشته باشد؛ خطای یک ردیف نباید دیتابیس را نیمه تغییر دهد.
\item GradeItem باید assignee داشته باشد و TA فقط آیتم assign شده را edit کند.
\item grade publish باید audit شود.
\item بعد از publish باید lock یا approval workflow برای تغییر نمره وجود داشته باشد.
\item Student فقط grade خودش را ببیند.
\item Professor و Head TA باید grade changes را ببینند.
\item regrade request فقط برای grade published قابل ثبت باشد.
\item برای هر grade record فقط یک regrade request باز مجاز باشد.
\item rubric برای نمره دهی دقیق اضافه شود.
\item assignment submission باید اضافه شود.
\item plagiarism/similarity check باید به عنوان قابلیت اختیاری یا placeholder اضافه شود.
\item late submission policy باید تعریف و enforce شود.
\end{enumerate}
\subsection{Messaging / Notification}
\begin{enumerate}[label=\arabic*. ,resume]
\item برای messaging باید security test روی XSS و unauthorized thread access نوشته شود.
\item unread count لازم است.
\item seen/read receipt لازم است.
\item message attachment لازم است.
\item close/reopen thread باید permission داشته باشد.
\item escalation از TA به Head TA یا Professor لازم است.
\item message rate limit لازم است.
\item notification باید channel داشته باشد: in-app، email، digest و urgent.
\item notification باید read/unread و mark all as read داشته باشد.
\item reminder برای deadlineها باید اضافه شود.
\end{enumerate}
\subsection{Office Hour / Sessions}
\begin{enumerate}[label=\arabic*. ,resume]
\item meeting link یا physical location، capacity، start/end، no double booking و ICS export باید با e2e تست شوند.
\item recurring sessions لازم است.
\item cancellation و reschedule لازم است.
\item waitlist برای capacity کامل لازم است.
\item attendance tracking لازم است.
\item feedback بعد از session لازم است.
\item automatic Meet creation از Google Calendar API می تواند در فاز بعدی اضافه شود.
\item host conflict باید در سطح DB/service تست شود.
\item student registration/cancel باید audit شود.
\item session reminder باید notification/email بفرستد.
\end{enumerate}
\section{لیست ۳۰۰ قابلیت موردنیاز پروژه همراه با مراحل پیاده سازی}
در این بخش، ۳۰۰ قابلیت پیشنهادی برای تبدیل TAFlow به یک سامانه کامل دانشگاهی آمده است. برای هر قابلیت، مسیر پیاده سازی و نوع تست لازم ذکر شده است.
\subsection{A) هویت، حساب کاربری و امنیت}

\begin{longtable}{|>{\raggedleft\arraybackslash}p{0.065\textwidth}|>{\raggedleft\arraybackslash}p{0.21\textwidth}|>{\raggedleft\arraybackslash}p{0.42\textwidth}|>{\raggedleft\arraybackslash}p{0.18\textwidth}|}
\hline
\textbf{شماره} & \textbf{قابلیت} & \textbf{مراحل پیاده سازی} & \textbf{تست و پذیرش} \\\hline
\endfirsthead
\hline
\textbf{شماره} & \textbf{قابلیت} & \textbf{مراحل پیاده سازی} & \textbf{تست و پذیرش} \\\hline
\endhead
1 & ثبت نام دانشجو با ایمیل دانشگاهی & مدل User/SecurityEvent و routeهای auth را تکمیل کن؛ validation، hashing/2FA/rate-limit را در service بگذار؛ UI تنظیمات حساب را بساز؛ unit/e2e برای login، status و audit بنویس. برای این قابلیت، ابتدا acceptance criteria بنویس، سپس migration/schema، service، API، UI، audit، test و documentation را در همان PR انجام بده. & Unit + integration + e2e؛ تست permission، validation، audit و حالت خطا. سناریوهای session، role، status و access denial را جداگانه تست کن. \\ \hline
2 & ورود با ایمیل و رمز & مدل User/SecurityEvent و routeهای auth را تکمیل کن؛ validation، hashing/2FA/rate-limit را در service بگذار؛ UI تنظیمات حساب را بساز؛ unit/e2e برای login، status و audit بنویس. برای این قابلیت، ابتدا acceptance criteria بنویس، سپس migration/schema، service، API، UI، audit، test و documentation را در همان PR انجام بده. & Unit + integration + e2e؛ تست permission، validation، audit و حالت خطا. سناریوهای session، role، status و access denial را جداگانه تست کن. \\ \hline
3 & ورود با SSO دانشگاهی / Keycloak & مدل User/SecurityEvent و routeهای auth را تکمیل کن؛ validation، hashing/2FA/rate-limit را در service بگذار؛ UI تنظیمات حساب را بساز؛ unit/e2e برای login، status و audit بنویس. برای این قابلیت، ابتدا acceptance criteria بنویس، سپس migration/schema، service، API، UI، audit، test و documentation را در همان PR انجام بده. & Unit + integration + e2e؛ تست permission، validation، audit و حالت خطا. سناریوهای session، role، status و access denial را جداگانه تست کن. \\ \hline
4 & تأیید ایمیل & مدل User/SecurityEvent و routeهای auth را تکمیل کن؛ validation، hashing/2FA/rate-limit را در service بگذار؛ UI تنظیمات حساب را بساز؛ unit/e2e برای login، status و audit بنویس. برای این قابلیت، ابتدا acceptance criteria بنویس، سپس migration/schema، service، API، UI، audit، test و documentation را در همان PR انجام بده. & Unit + integration + e2e؛ تست permission، validation، audit و حالت خطا. سناریوهای session، role، status و access denial را جداگانه تست کن. \\ \hline
5 & فراموشی رمز عبور & مدل User/SecurityEvent و routeهای auth را تکمیل کن؛ validation، hashing/2FA/rate-limit را در service بگذار؛ UI تنظیمات حساب را بساز؛ unit/e2e برای login، status و audit بنویس. برای این قابلیت، ابتدا acceptance criteria بنویس، سپس migration/schema، service، API، UI، audit، test و documentation را در همان PR انجام بده. & Unit + integration + e2e؛ تست permission، validation، audit و حالت خطا. سناریوهای session، role، status و access denial را جداگانه تست کن. \\ \hline
6 & تغییر رمز عبور & مدل User/SecurityEvent و routeهای auth را تکمیل کن؛ validation، hashing/2FA/rate-limit را در service بگذار؛ UI تنظیمات حساب را بساز؛ unit/e2e برای login، status و audit بنویس. برای این قابلیت، ابتدا acceptance criteria بنویس، سپس migration/schema، service، API، UI، audit، test و documentation را در همان PR انجام بده. & Unit + integration + e2e؛ تست permission، validation، audit و حالت خطا. سناریوهای session، role، status و access denial را جداگانه تست کن. \\ \hline
7 & اجبار تغییر رمز بعد از reset & مدل User/SecurityEvent و routeهای auth را تکمیل کن؛ validation، hashing/2FA/rate-limit را در service بگذار؛ UI تنظیمات حساب را بساز؛ unit/e2e برای login، status و audit بنویس. برای این قابلیت، ابتدا acceptance criteria بنویس، سپس migration/schema، service، API، UI، audit، test و documentation را در همان PR انجام بده. & Unit + integration + e2e؛ تست permission، validation، audit و حالت خطا. سناریوهای session، role، status و access denial را جداگانه تست کن. \\ \hline
8 & احراز هویت دو مرحله ای TOTP & مدل User/SecurityEvent و routeهای auth را تکمیل کن؛ validation، hashing/2FA/rate-limit را در service بگذار؛ UI تنظیمات حساب را بساز؛ unit/e2e برای login، status و audit بنویس. برای این قابلیت، ابتدا acceptance criteria بنویس، سپس migration/schema، service، API، UI، audit، test و documentation را در همان PR انجام بده. & Unit + integration + e2e؛ تست permission، validation، audit و حالت خطا. سناریوهای session، role، status و access denial را جداگانه تست کن. \\ \hline
9 & کدهای بازیابی 2FA & مدل User/SecurityEvent و routeهای auth را تکمیل کن؛ validation، hashing/2FA/rate-limit را در service بگذار؛ UI تنظیمات حساب را بساز؛ unit/e2e برای login، status و audit بنویس. برای این قابلیت، ابتدا acceptance criteria بنویس، سپس migration/schema، service، API، UI، audit، test و documentation را در همان PR انجام بده. & Unit + integration + e2e؛ تست permission، validation، audit و حالت خطا. سناریوهای session، role، status و access denial را جداگانه تست کن. \\ \hline
10 & قفل موقت حساب بعد از تلاش ناموفق & مدل User/SecurityEvent و routeهای auth را تکمیل کن؛ validation، hashing/2FA/rate-limit را در service بگذار؛ UI تنظیمات حساب را بساز؛ unit/e2e برای login، status و audit بنویس. برای این قابلیت، ابتدا acceptance criteria بنویس، سپس migration/schema، service، API، UI، audit، test و documentation را در همان PR انجام بده. & Unit + integration + e2e؛ تست permission، validation، audit و حالت خطا. سناریوهای session، role، status و access denial را جداگانه تست کن. \\ \hline
11 & rate limit برای login/register & مدل User/SecurityEvent و routeهای auth را تکمیل کن؛ validation، hashing/2FA/rate-limit را در service بگذار؛ UI تنظیمات حساب را بساز؛ unit/e2e برای login، status و audit بنویس. برای این قابلیت، ابتدا acceptance criteria بنویس، سپس migration/schema، service، API، UI، audit، test و documentation را در همان PR انجام بده. & Unit + integration + e2e؛ تست permission، validation، audit و حالت خطا. سناریوهای session، role، status و access denial را جداگانه تست کن. \\ \hline
12 & مدیریت sessionهای فعال & مدل User/SecurityEvent و routeهای auth را تکمیل کن؛ validation، hashing/2FA/rate-limit را در service بگذار؛ UI تنظیمات حساب را بساز؛ unit/e2e برای login، status و audit بنویس. برای این قابلیت، ابتدا acceptance criteria بنویس، سپس migration/schema، service، API، UI، audit، test و documentation را در همان PR انجام بده. & Unit + integration + e2e؛ تست permission، validation، audit و حالت خطا. سناریوهای session، role، status و access denial را جداگانه تست کن. \\ \hline
13 & خروج از همه دستگاه ها & مدل User/SecurityEvent و routeهای auth را تکمیل کن؛ validation، hashing/2FA/rate-limit را در service بگذار؛ UI تنظیمات حساب را بساز؛ unit/e2e برای login، status و audit بنویس. برای این قابلیت، ابتدا acceptance criteria بنویس، سپس migration/schema، service، API، UI، audit، test و documentation را در همان PR انجام بده. & Unit + integration + e2e؛ تست permission، validation، audit و حالت خطا. سناریوهای session، role، status و access denial را جداگانه تست کن. \\ \hline
14 & وضعیت حساب: active/suspended/deleted & مدل User/SecurityEvent و routeهای auth را تکمیل کن؛ validation، hashing/2FA/rate-limit را در service بگذار؛ UI تنظیمات حساب را بساز؛ unit/e2e برای login، status و audit بنویس. برای این قابلیت، ابتدا acceptance criteria بنویس، سپس migration/schema، service، API، UI، audit، test و documentation را در همان PR انجام بده. & Unit + integration + e2e؛ تست permission، validation، audit و حالت خطا. سناریوهای session، role، status و access denial را جداگانه تست کن. \\ \hline
15 & پروفایل دانشجو & مدل User/SecurityEvent و routeهای auth را تکمیل کن؛ validation، hashing/2FA/rate-limit را در service بگذار؛ UI تنظیمات حساب را بساز؛ unit/e2e برای login، status و audit بنویس. برای این قابلیت، ابتدا acceptance criteria بنویس، سپس migration/schema، service، API، UI، audit، test و documentation را در همان PR انجام بده. & Unit + integration + e2e؛ تست permission، validation، audit و حالت خطا. سناریوهای session، role، status و access denial را جداگانه تست کن. \\ \hline
16 & پروفایل استاد & مدل User/SecurityEvent و routeهای auth را تکمیل کن؛ validation، hashing/2FA/rate-limit را در service بگذار؛ UI تنظیمات حساب را بساز؛ unit/e2e برای login، status و audit بنویس. برای این قابلیت، ابتدا acceptance criteria بنویس، سپس migration/schema، service، API، UI، audit، test و documentation را در همان PR انجام بده. & Unit + integration + e2e؛ تست permission، validation، audit و حالت خطا. سناریوهای session، role، status و access denial را جداگانه تست کن. \\ \hline
17 & آپلود عکس پروفایل & مدل User/SecurityEvent و routeهای auth را تکمیل کن؛ validation، hashing/2FA/rate-limit را در service بگذار؛ UI تنظیمات حساب را بساز؛ unit/e2e برای login، status و audit بنویس. برای این قابلیت، ابتدا acceptance criteria بنویس، سپس migration/schema، service، API، UI، audit، test و documentation را در همان PR انجام بده. & Unit + integration + e2e؛ تست permission، validation، audit و حالت خطا. سناریوهای session، role، status و access denial را جداگانه تست کن. \\ \hline
18 & مدیریت timezone و زبان & مدل User/SecurityEvent و routeهای auth را تکمیل کن؛ validation، hashing/2FA/rate-limit را در service بگذار؛ UI تنظیمات حساب را بساز؛ unit/e2e برای login، status و audit بنویس. برای این قابلیت، ابتدا acceptance criteria بنویس، سپس migration/schema، service، API، UI، audit، test و documentation را در همان PR انجام بده. & Unit + integration + e2e؛ تست permission، validation، audit و حالت خطا. سناریوهای session، role، status و access denial را جداگانه تست کن. \\ \hline
19 & تنظیمات اعلان ها & مدل User/SecurityEvent و routeهای auth را تکمیل کن؛ validation، hashing/2FA/rate-limit را در service بگذار؛ UI تنظیمات حساب را بساز؛ unit/e2e برای login، status و audit بنویس. برای این قابلیت، ابتدا acceptance criteria بنویس، سپس migration/schema، service، API، UI، audit، test و documentation را در همان PR انجام بده. & Unit + integration + e2e؛ تست permission، validation، audit و حالت خطا. سناریوهای session، role، status و access denial را جداگانه تست کن. \\ \hline
20 & audit log برای عملیات امنیتی & مدل User/SecurityEvent و routeهای auth را تکمیل کن؛ validation، hashing/2FA/rate-limit را در service بگذار؛ UI تنظیمات حساب را بساز؛ unit/e2e برای login، status و audit بنویس. برای این قابلیت، ابتدا acceptance criteria بنویس، سپس migration/schema، service، API، UI، audit، test و documentation را در همان PR انجام بده. & Unit + integration + e2e؛ تست permission، validation، audit و حالت خطا. سناریوهای session، role، status و access denial را جداگانه تست کن. \\ \hline
\end{longtable}
\subsection{B) ساختار دانشگاهی}

\begin{longtable}{|>{\raggedleft\arraybackslash}p{0.065\textwidth}|>{\raggedleft\arraybackslash}p{0.21\textwidth}|>{\raggedleft\arraybackslash}p{0.42\textwidth}|>{\raggedleft\arraybackslash}p{0.18\textwidth}|}
\hline
\textbf{شماره} & \textbf{قابلیت} & \textbf{مراحل پیاده سازی} & \textbf{تست و پذیرش} \\\hline
\endfirsthead
\hline
\textbf{شماره} & \textbf{قابلیت} & \textbf{مراحل پیاده سازی} & \textbf{تست و پذیرش} \\\hline
\endhead
21 & مدیریت دانشکده/دپارتمان & مدل های academic را در Prisma و migration تثبیت کن؛ import/export و validation ظرفیت/ترم را در service اضافه کن؛ صفحه admin و تست seed/migration/CSV بنویس. برای این قابلیت، ابتدا acceptance criteria بنویس، سپس migration/schema، service، API، UI، audit، test و documentation را در همان PR انجام بده. & Unit + integration + e2e؛ تست permission، validation، audit و حالت خطا. \\ \hline
22 & مدیریت ترم ها & مدل های academic را در Prisma و migration تثبیت کن؛ import/export و validation ظرفیت/ترم را در service اضافه کن؛ صفحه admin و تست seed/migration/CSV بنویس. برای این قابلیت، ابتدا acceptance criteria بنویس، سپس migration/schema، service، API، UI، audit، test و documentation را در همان PR انجام بده. & Unit + integration + e2e؛ تست permission، validation، audit و حالت خطا. \\ \hline
23 & مدیریت درس ها & مدل های academic را در Prisma و migration تثبیت کن؛ import/export و validation ظرفیت/ترم را در service اضافه کن؛ صفحه admin و تست seed/migration/CSV بنویس. برای این قابلیت، ابتدا acceptance criteria بنویس، سپس migration/schema، service، API، UI، audit، test و documentation را در همان PR انجام بده. & Unit + integration + e2e؛ تست permission، validation، audit و حالت خطا. \\ \hline
24 & مدیریت ارائه درس در هر ترم & مدل های academic را در Prisma و migration تثبیت کن؛ import/export و validation ظرفیت/ترم را در service اضافه کن؛ صفحه admin و تست seed/migration/CSV بنویس. برای این قابلیت، ابتدا acceptance criteria بنویس، سپس migration/schema، service، API، UI، audit، test و documentation را در همان PR انجام بده. & Unit + integration + e2e؛ تست permission، validation، audit و حالت خطا. \\ \hline
25 & تفکیک Course از CourseOffering & مدل های academic را در Prisma و migration تثبیت کن؛ import/export و validation ظرفیت/ترم را در service اضافه کن؛ صفحه admin و تست seed/migration/CSV بنویس. برای این قابلیت، ابتدا acceptance criteria بنویس، سپس migration/schema، service، API، UI، audit، test و documentation را در همان PR انجام بده. & Unit + integration + e2e؛ تست permission، validation، audit و حالت خطا. \\ \hline
26 & مدیریت section/class number & مدل های academic را در Prisma و migration تثبیت کن؛ import/export و validation ظرفیت/ترم را در service اضافه کن؛ صفحه admin و تست seed/migration/CSV بنویس. برای این قابلیت، ابتدا acceptance criteria بنویس، سپس migration/schema، service، API، UI، audit، test و documentation را در همان PR انجام بده. & Unit + integration + e2e؛ تست permission، validation، audit و حالت خطا. \\ \hline
27 & مدیریت ظرفیت کلاس & مدل های academic را در Prisma و migration تثبیت کن؛ import/export و validation ظرفیت/ترم را در service اضافه کن؛ صفحه admin و تست seed/migration/CSV بنویس. برای این قابلیت، ابتدا acceptance criteria بنویس، سپس migration/schema، service، API، UI، audit، test و documentation را در همان PR انجام بده. & Unit + integration + e2e؛ تست permission، validation، audit و حالت خطا. \\ \hline
28 & import دروس از Excel & مدل های academic را در Prisma و migration تثبیت کن؛ import/export و validation ظرفیت/ترم را در service اضافه کن؛ صفحه admin و تست seed/migration/CSV بنویس. برای این قابلیت، ابتدا acceptance criteria بنویس، سپس migration/schema، service، API، UI، audit، test و documentation را در همان PR انجام بده. & Unit + integration + e2e؛ تست permission، validation، audit و حالت خطا. \\ \hline
29 & import دانشجوها از Excel & مدل های academic را در Prisma و migration تثبیت کن؛ import/export و validation ظرفیت/ترم را در service اضافه کن؛ صفحه admin و تست seed/migration/CSV بنویس. برای این قابلیت، ابتدا acceptance criteria بنویس، سپس migration/schema، service، API، UI، audit، test و documentation را در همان PR انجام بده. & Unit + integration + e2e؛ تست permission، validation، audit و حالت خطا. \\ \hline
30 & import استادها از Excel & مدل های academic را در Prisma و migration تثبیت کن؛ import/export و validation ظرفیت/ترم را در service اضافه کن؛ صفحه admin و تست seed/migration/CSV بنویس. برای این قابلیت، ابتدا acceptance criteria بنویس، سپس migration/schema، service، API، UI، audit، test و documentation را در همان PR انجام بده. & Unit + integration + e2e؛ تست permission، validation، audit و حالت خطا. \\ \hline
31 & ثبت نام دانشجو در course offering & مدل های academic را در Prisma و migration تثبیت کن؛ import/export و validation ظرفیت/ترم را در service اضافه کن؛ صفحه admin و تست seed/migration/CSV بنویس. برای این قابلیت، ابتدا acceptance criteria بنویس، سپس migration/schema، service، API، UI، audit، test و documentation را در همان PR انجام بده. & Unit + integration + e2e؛ تست permission، validation، audit و حالت خطا. \\ \hline
32 & drop/re-enroll دانشجو & مدل های academic را در Prisma و migration تثبیت کن؛ import/export و validation ظرفیت/ترم را در service اضافه کن؛ صفحه admin و تست seed/migration/CSV بنویس. برای این قابلیت، ابتدا acceptance criteria بنویس، سپس migration/schema، service، API، UI، audit، test و documentation را در همان PR انجام بده. & Unit + integration + e2e؛ تست permission، validation، audit و حالت خطا. \\ \hline
33 & مشاهده roster کلاس & مدل های academic را در Prisma و migration تثبیت کن؛ import/export و validation ظرفیت/ترم را در service اضافه کن؛ صفحه admin و تست seed/migration/CSV بنویس. برای این قابلیت، ابتدا acceptance criteria بنویس، سپس migration/schema، service، API، UI، audit، test و documentation را در همان PR انجام بده. & Unit + integration + e2e؛ تست permission، validation، audit و حالت خطا. \\ \hline
34 & export roster به CSV/XLSX & مدل های academic را در Prisma و migration تثبیت کن؛ import/export و validation ظرفیت/ترم را در service اضافه کن؛ صفحه admin و تست seed/migration/CSV بنویس. برای این قابلیت، ابتدا acceptance criteria بنویس، سپس migration/schema، service، API، UI، audit، test و documentation را در همان PR انجام بده. & Unit + integration + e2e؛ تست permission، validation، audit و حالت خطا. \\ \hline
35 & آرشیو کردن ترم های قدیمی & مدل های academic را در Prisma و migration تثبیت کن؛ import/export و validation ظرفیت/ترم را در service اضافه کن؛ صفحه admin و تست seed/migration/CSV بنویس. برای این قابلیت، ابتدا acceptance criteria بنویس، سپس migration/schema، service، API، UI، audit، test و documentation را در همان PR انجام بده. & Unit + integration + e2e؛ تست permission، validation، audit و حالت خطا. \\ \hline
\end{longtable}
\subsection{C) نقش ها و دسترسی ها}

\begin{longtable}{|>{\raggedleft\arraybackslash}p{0.065\textwidth}|>{\raggedleft\arraybackslash}p{0.21\textwidth}|>{\raggedleft\arraybackslash}p{0.42\textwidth}|>{\raggedleft\arraybackslash}p{0.18\textwidth}|}
\hline
\textbf{شماره} & \textbf{قابلیت} & \textbf{مراحل پیاده سازی} & \textbf{تست و پذیرش} \\\hline
\endfirsthead
\hline
\textbf{شماره} & \textbf{قابلیت} & \textbf{مراحل پیاده سازی} & \textbf{تست و پذیرش} \\\hline
\endhead
36 & نقش global: student/professor/admin & Permission را server-side و course-scoped کن؛ guard مشترک API بساز؛ UI را permission-aware کن؛ تست escalation، revoked و expired role بنویس. برای این قابلیت، ابتدا acceptance criteria بنویس، سپس migration/schema، service، API، UI، audit، test و documentation را در همان PR انجام بده. & Unit + integration + e2e؛ تست permission، validation، audit و حالت خطا. سناریوهای session، role، status و access denial را جداگانه تست کن. \\ \hline
37 & نقش درس محور: student/TA/Head TA/professor & Permission را server-side و course-scoped کن؛ guard مشترک API بساز؛ UI را permission-aware کن؛ تست escalation، revoked و expired role بنویس. برای این قابلیت، ابتدا acceptance criteria بنویس، سپس migration/schema، service، API، UI، audit، test و documentation را در همان PR انجام بده. & Unit + integration + e2e؛ تست permission، validation، audit و حالت خطا. سناریوهای session، role، status و access denial را جداگانه تست کن. \\ \hline
38 & assign role برای یک course offering & Permission را server-side و course-scoped کن؛ guard مشترک API بساز؛ UI را permission-aware کن؛ تست escalation، revoked و expired role بنویس. برای این قابلیت، ابتدا acceptance criteria بنویس، سپس migration/schema، service، API، UI، audit، test و documentation را در همان PR انجام بده. & Unit + integration + e2e؛ تست permission، validation، audit و حالت خطا. سناریوهای session، role، status و access denial را جداگانه تست کن. \\ \hline
39 & revoke role & Permission را server-side و course-scoped کن؛ guard مشترک API بساز؛ UI را permission-aware کن؛ تست escalation، revoked و expired role بنویس. برای این قابلیت، ابتدا acceptance criteria بنویس، سپس migration/schema، service، API، UI، audit، test و documentation را در همان PR انجام بده. & Unit + integration + e2e؛ تست permission، validation، audit و حالت خطا. سناریوهای session، role، status و access denial را جداگانه تست کن. \\ \hline
40 & role expiration & Permission را server-side و course-scoped کن؛ guard مشترک API بساز؛ UI را permission-aware کن؛ تست escalation، revoked و expired role بنویس. برای این قابلیت، ابتدا acceptance criteria بنویس، سپس migration/schema، service، API، UI، audit، test و documentation را در همان PR انجام بده. & Unit + integration + e2e؛ تست permission، validation، audit و حالت خطا. سناریوهای session، role، status و access denial را جداگانه تست کن. \\ \hline
41 & role history & Permission را server-side و course-scoped کن؛ guard مشترک API بساز؛ UI را permission-aware کن؛ تست escalation، revoked و expired role بنویس. برای این قابلیت، ابتدا acceptance criteria بنویس، سپس migration/schema، service، API، UI، audit، test و documentation را در همان PR انجام بده. & Unit + integration + e2e؛ تست permission، validation، audit و حالت خطا. سناریوهای session، role، status و access denial را جداگانه تست کن. \\ \hline
42 & permission override با permissionsJson & Permission را server-side و course-scoped کن؛ guard مشترک API بساز؛ UI را permission-aware کن؛ تست escalation، revoked و expired role بنویس. برای این قابلیت، ابتدا acceptance criteria بنویس، سپس migration/schema، service، API، UI، audit، test و documentation را در همان PR انجام بده. & Unit + integration + e2e؛ تست permission، validation، audit و حالت خطا. سناریوهای session، role، status و access denial را جداگانه تست کن. \\ \hline
43 & Head TA management permissions & Permission را server-side و course-scoped کن؛ guard مشترک API بساز؛ UI را permission-aware کن؛ تست escalation، revoked و expired role بنویس. برای این قابلیت، ابتدا acceptance criteria بنویس، سپس migration/schema، service، API، UI، audit، test و documentation را در همان PR انجام بده. & Unit + integration + e2e؛ تست permission، validation، audit و حالت خطا. سناریوهای session، role، status و access denial را جداگانه تست کن. \\ \hline
44 & Professor-only permissions & Permission را server-side و course-scoped کن؛ guard مشترک API بساز؛ UI را permission-aware کن؛ تست escalation، revoked و expired role بنویس. برای این قابلیت، ابتدا acceptance criteria بنویس، سپس migration/schema، service، API، UI، audit، test و documentation را در همان PR انجام بده. & Unit + integration + e2e؛ تست permission، validation، audit و حالت خطا. سناریوهای session، role، status و access denial را جداگانه تست کن. \\ \hline
45 & Education Admin permissions & Permission را server-side و course-scoped کن؛ guard مشترک API بساز؛ UI را permission-aware کن؛ تست escalation، revoked و expired role بنویس. برای این قابلیت، ابتدا acceptance criteria بنویس، سپس migration/schema، service، API، UI، audit، test و documentation را در همان PR انجام بده. & Unit + integration + e2e؛ تست permission، validation، audit و حالت خطا. سناریوهای session، role، status و access denial را جداگانه تست کن. \\ \hline
46 & permission-aware search & Permission را server-side و course-scoped کن؛ guard مشترک API بساز؛ UI را permission-aware کن؛ تست escalation، revoked و expired role بنویس. برای این قابلیت، ابتدا acceptance criteria بنویس، سپس migration/schema، service، API، UI، audit، test و documentation را در همان PR انجام بده. & Unit + integration + e2e؛ تست permission، validation، audit و حالت خطا. سناریوهای session، role، status و access denial را جداگانه تست کن. \\ \hline
47 & permission-aware dashboards & Permission را server-side و course-scoped کن؛ guard مشترک API بساز؛ UI را permission-aware کن؛ تست escalation، revoked و expired role بنویس. برای این قابلیت، ابتدا acceptance criteria بنویس، سپس migration/schema، service، API، UI، audit، test و documentation را در همان PR انجام بده. & Unit + integration + e2e؛ تست permission، validation، audit و حالت خطا. سناریوهای session، role، status و access denial را جداگانه تست کن. \\ \hline
48 & permission audit & Permission را server-side و course-scoped کن؛ guard مشترک API بساز؛ UI را permission-aware کن؛ تست escalation، revoked و expired role بنویس. برای این قابلیت، ابتدا acceptance criteria بنویس، سپس migration/schema، service، API، UI، audit، test و documentation را در همان PR انجام بده. & Unit + integration + e2e؛ تست permission، validation، audit و حالت خطا. سناریوهای session، role، status و access denial را جداگانه تست کن. \\ \hline
49 & تست escalation & Permission را server-side و course-scoped کن؛ guard مشترک API بساز؛ UI را permission-aware کن؛ تست escalation، revoked و expired role بنویس. برای این قابلیت، ابتدا acceptance criteria بنویس، سپس migration/schema، service، API، UI، audit، test و documentation را در همان PR انجام بده. & Unit + integration + e2e؛ تست permission، validation، audit و حالت خطا. سناریوهای session، role، status و access denial را جداگانه تست کن. \\ \hline
50 & تست revoked/expired role & Permission را server-side و course-scoped کن؛ guard مشترک API بساز؛ UI را permission-aware کن؛ تست escalation، revoked و expired role بنویس. برای این قابلیت، ابتدا acceptance criteria بنویس، سپس migration/schema، service، API، UI، audit، test و documentation را در همان PR انجام بده. & Unit + integration + e2e؛ تست permission، validation، audit و حالت خطا. سناریوهای session، role، status و access denial را جداگانه تست کن. \\ \hline
\end{longtable}
\subsection{D) درخواست و انتخاب TA/Head TA}

\begin{longtable}{|>{\raggedleft\arraybackslash}p{0.065\textwidth}|>{\raggedleft\arraybackslash}p{0.21\textwidth}|>{\raggedleft\arraybackslash}p{0.42\textwidth}|>{\raggedleft\arraybackslash}p{0.18\textwidth}|}
\hline
\textbf{شماره} & \textbf{قابلیت} & \textbf{مراحل پیاده سازی} & \textbf{تست و پذیرش} \\\hline
\endfirsthead
\hline
\textbf{شماره} & \textbf{قابلیت} & \textbf{مراحل پیاده سازی} & \textbf{تست و پذیرش} \\\hline
\endhead
51 & ساخت فرصت TA توسط استاد & workflow وضعیت application را با transaction پیاده کن؛ API status/interview/withdraw را تکمیل کن؛ notification/audit را اجباری کن؛ e2e مسیر student-professor-headTA بنویس. برای این قابلیت، ابتدا acceptance criteria بنویس، سپس migration/schema، service، API، UI، audit، test و documentation را در همان PR انجام بده. & Unit + integration + e2e؛ تست permission، validation، audit و حالت خطا. \\ \hline
52 & publish/close/cancel فرصت TA & workflow وضعیت application را با transaction پیاده کن؛ API status/interview/withdraw را تکمیل کن؛ notification/audit را اجباری کن؛ e2e مسیر student-professor-headTA بنویس. برای این قابلیت، ابتدا acceptance criteria بنویس، سپس migration/schema، service، API، UI، audit، test و documentation را در همان PR انجام بده. & Unit + integration + e2e؛ تست permission، validation، audit و حالت خطا. \\ \hline
53 & تعیین deadline & workflow وضعیت application را با transaction پیاده کن؛ API status/interview/withdraw را تکمیل کن؛ notification/audit را اجباری کن؛ e2e مسیر student-professor-headTA بنویس. برای این قابلیت، ابتدا acceptance criteria بنویس، سپس migration/schema، service، API، UI، audit، test و documentation را در همان PR انجام بده. & Unit + integration + e2e؛ تست permission، validation، audit و حالت خطا. \\ \hline
54 & تعیین تعداد TA مورد نیاز & workflow وضعیت application را با transaction پیاده کن؛ API status/interview/withdraw را تکمیل کن؛ notification/audit را اجباری کن؛ e2e مسیر student-professor-headTA بنویس. برای این قابلیت، ابتدا acceptance criteria بنویس، سپس migration/schema، service، API، UI، audit، test و documentation را در همان PR انجام بده. & Unit + integration + e2e؛ تست permission، validation، audit و حالت خطا. \\ \hline
55 & تعیین نیاز به Head TA & workflow وضعیت application را با transaction پیاده کن؛ API status/interview/withdraw را تکمیل کن؛ notification/audit را اجباری کن؛ e2e مسیر student-professor-headTA بنویس. برای این قابلیت، ابتدا acceptance criteria بنویس، سپس migration/schema، service، API، UI، audit، test و documentation را در همان PR انجام بده. & Unit + integration + e2e؛ تست permission، validation، audit و حالت خطا. \\ \hline
56 & تعیین rubric انتخاب & workflow وضعیت application را با transaction پیاده کن؛ API status/interview/withdraw را تکمیل کن؛ notification/audit را اجباری کن؛ e2e مسیر student-professor-headTA بنویس. برای این قابلیت، ابتدا acceptance criteria بنویس، سپس migration/schema، service، API، UI، audit، test و documentation را در همان PR انجام بده. & Unit + integration + e2e؛ تست permission، validation، audit و حالت خطا. \\ \hline
57 & فرم درخواست دانشجو & workflow وضعیت application را با transaction پیاده کن؛ API status/interview/withdraw را تکمیل کن؛ notification/audit را اجباری کن؛ e2e مسیر student-professor-headTA بنویس. برای این قابلیت، ابتدا acceptance criteria بنویس، سپس migration/schema، service، API، UI، audit، test و documentation را در همان PR انجام بده. & Unit + integration + e2e؛ تست permission، validation، audit و حالت خطا. \\ \hline
58 & انتخاب نقش درخواستی: TA/Head TA/Either & workflow وضعیت application را با transaction پیاده کن؛ API status/interview/withdraw را تکمیل کن؛ notification/audit را اجباری کن؛ e2e مسیر student-professor-headTA بنویس. برای این قابلیت، ابتدا acceptance criteria بنویس، سپس migration/schema، service، API، UI، audit، test و documentation را در همان PR انجام بده. & Unit + integration + e2e؛ تست permission، validation، audit و حالت خطا. \\ \hline
59 & آپلود رزومه & workflow وضعیت application را با transaction پیاده کن؛ API status/interview/withdraw را تکمیل کن؛ notification/audit را اجباری کن؛ e2e مسیر student-professor-headTA بنویس. برای این قابلیت، ابتدا acceptance criteria بنویس، سپس migration/schema، service، API، UI، audit، test و documentation را در همان PR انجام بده. & Unit + integration + e2e؛ تست permission، validation، audit و حالت خطا. \\ \hline
60 & انگیزه نامه & workflow وضعیت application را با transaction پیاده کن؛ API status/interview/withdraw را تکمیل کن؛ notification/audit را اجباری کن؛ e2e مسیر student-professor-headTA بنویس. برای این قابلیت، ابتدا acceptance criteria بنویس، سپس migration/schema، service، API، UI، audit، test و documentation را در همان PR انجام بده. & Unit + integration + e2e؛ تست permission، validation، audit و حالت خطا. \\ \hline
61 & جلوگیری از double apply & workflow وضعیت application را با transaction پیاده کن؛ API status/interview/withdraw را تکمیل کن؛ notification/audit را اجباری کن؛ e2e مسیر student-professor-headTA بنویس. برای این قابلیت، ابتدا acceptance criteria بنویس، سپس migration/schema، service، API، UI، audit، test و documentation را در همان PR انجام بده. & Unit + integration + e2e؛ تست permission، validation، audit و حالت خطا. \\ \hline
62 & بررسی application توسط استاد & workflow وضعیت application را با transaction پیاده کن؛ API status/interview/withdraw را تکمیل کن؛ notification/audit را اجباری کن؛ e2e مسیر student-professor-headTA بنویس. برای این قابلیت، ابتدا acceptance criteria بنویس، سپس migration/schema، service، API، UI، audit، test و documentation را در همان PR انجام بده. & Unit + integration + e2e؛ تست permission، validation، audit و حالت خطا. \\ \hline
63 & بررسی application توسط Head TA & workflow وضعیت application را با transaction پیاده کن؛ API status/interview/withdraw را تکمیل کن؛ notification/audit را اجباری کن؛ e2e مسیر student-professor-headTA بنویس. برای این قابلیت، ابتدا acceptance criteria بنویس، سپس migration/schema، service، API، UI، audit، test و documentation را در همان PR انجام بده. & Unit + integration + e2e؛ تست permission، validation، audit و حالت خطا. \\ \hline
64 & shortlist کردن & workflow وضعیت application را با transaction پیاده کن؛ API status/interview/withdraw را تکمیل کن؛ notification/audit را اجباری کن؛ e2e مسیر student-professor-headTA بنویس. برای این قابلیت، ابتدا acceptance criteria بنویس، سپس migration/schema، service، API، UI، audit، test و documentation را در همان PR انجام بده. & Unit + integration + e2e؛ تست permission، validation، audit و حالت خطا. \\ \hline
65 & دعوت به مصاحبه & workflow وضعیت application را با transaction پیاده کن؛ API status/interview/withdraw را تکمیل کن؛ notification/audit را اجباری کن؛ e2e مسیر student-professor-headTA بنویس. برای این قابلیت، ابتدا acceptance criteria بنویس، سپس migration/schema، service، API، UI، audit، test و documentation را در همان PR انجام بده. & Unit + integration + e2e؛ تست permission، validation، audit و حالت خطا. \\ \hline
66 & ثبت زمان مصاحبه & workflow وضعیت application را با transaction پیاده کن؛ API status/interview/withdraw را تکمیل کن؛ notification/audit را اجباری کن؛ e2e مسیر student-professor-headTA بنویس. برای این قابلیت، ابتدا acceptance criteria بنویس، سپس migration/schema، service، API، UI، audit، test و documentation را در همان PR انجام بده. & Unit + integration + e2e؛ تست permission، validation، audit و حالت خطا. \\ \hline
67 & ثبت نمره مصاحبه & workflow وضعیت application را با transaction پیاده کن؛ API status/interview/withdraw را تکمیل کن؛ notification/audit را اجباری کن؛ e2e مسیر student-professor-headTA بنویس. برای این قابلیت، ابتدا acceptance criteria بنویس، سپس migration/schema، service، API، UI، audit، test و documentation را در همان PR انجام بده. & Unit + integration + e2e؛ تست permission، validation، audit و حالت خطا. \\ \hline
68 & ثبت note داخلی & workflow وضعیت application را با transaction پیاده کن؛ API status/interview/withdraw را تکمیل کن؛ notification/audit را اجباری کن؛ e2e مسیر student-professor-headTA بنویس. برای این قابلیت، ابتدا acceptance criteria بنویس، سپس migration/schema، service، API، UI، audit، test و documentation را در همان PR انجام بده. & Unit + integration + e2e؛ تست permission، validation، audit و حالت خطا. \\ \hline
69 & accept/reject نهایی & workflow وضعیت application را با transaction پیاده کن؛ API status/interview/withdraw را تکمیل کن؛ notification/audit را اجباری کن؛ e2e مسیر student-professor-headTA بنویس. برای این قابلیت، ابتدا acceptance criteria بنویس، سپس migration/schema، service، API، UI، audit، test و documentation را در همان PR انجام بده. & Unit + integration + e2e؛ تست permission، validation، audit و حالت خطا. \\ \hline
70 & ساخت خودکار CourseRoleAssignment بعد از accept & workflow وضعیت application را با transaction پیاده کن؛ API status/interview/withdraw را تکمیل کن؛ notification/audit را اجباری کن؛ e2e مسیر student-professor-headTA بنویس. برای این قابلیت، ابتدا acceptance criteria بنویس، سپس migration/schema، service، API، UI، audit، test و documentation را در همان PR انجام بده. & Unit + integration + e2e؛ تست permission، validation، audit و حالت خطا. \\ \hline
71 & ارسال notification بعد از تغییر status & workflow وضعیت application را با transaction پیاده کن؛ API status/interview/withdraw را تکمیل کن؛ notification/audit را اجباری کن؛ e2e مسیر student-professor-headTA بنویس. برای این قابلیت، ابتدا acceptance criteria بنویس، سپس migration/schema، service، API، UI، audit، test و documentation را در همان PR انجام بده. & Unit + integration + e2e؛ تست permission، validation، audit و حالت خطا. \\ \hline
72 & ثبت audit برای هر تغییر status & workflow وضعیت application را با transaction پیاده کن؛ API status/interview/withdraw را تکمیل کن؛ notification/audit را اجباری کن؛ e2e مسیر student-professor-headTA بنویس. برای این قابلیت، ابتدا acceptance criteria بنویس، سپس migration/schema، service، API، UI، audit، test و documentation را در همان PR انجام بده. & Unit + integration + e2e؛ تست permission، validation، audit و حالت خطا. \\ \hline
73 & امکان withdraw توسط دانشجو & workflow وضعیت application را با transaction پیاده کن؛ API status/interview/withdraw را تکمیل کن؛ notification/audit را اجباری کن؛ e2e مسیر student-professor-headTA بنویس. برای این قابلیت، ابتدا acceptance criteria بنویس، سپس migration/schema، service، API، UI، audit، test و documentation را در همان PR انجام بده. & Unit + integration + e2e؛ تست permission، validation، audit و حالت خطا. \\ \hline
74 & appeal/review بعد از rejection & workflow وضعیت application را با transaction پیاده کن؛ API status/interview/withdraw را تکمیل کن؛ notification/audit را اجباری کن؛ e2e مسیر student-professor-headTA بنویس. برای این قابلیت، ابتدا acceptance criteria بنویس، سپس migration/schema، service، API، UI، audit، test و documentation را در همان PR انجام بده. & Unit + integration + e2e؛ تست permission، validation، audit و حالت خطا. \\ \hline
75 & recommendation engine برای پیشنهاد TA & workflow وضعیت application را با transaction پیاده کن؛ API status/interview/withdraw را تکمیل کن؛ notification/audit را اجباری کن؛ e2e مسیر student-professor-headTA بنویس. برای این قابلیت، ابتدا acceptance criteria بنویس، سپس migration/schema، service، API، UI، audit، test و documentation را در همان PR انجام بده. & Unit + integration + e2e؛ تست permission، validation، audit و حالت خطا. \\ \hline
\end{longtable}
\subsection{E) Talent Pool و رزومه آموزشی}

\begin{longtable}{|>{\raggedleft\arraybackslash}p{0.065\textwidth}|>{\raggedleft\arraybackslash}p{0.21\textwidth}|>{\raggedleft\arraybackslash}p{0.42\textwidth}|>{\raggedleft\arraybackslash}p{0.18\textwidth}|}
\hline
\textbf{شماره} & \textbf{قابلیت} & \textbf{مراحل پیاده سازی} & \textbf{تست و پذیرش} \\\hline
\endfirsthead
\hline
\textbf{شماره} & \textbf{قابلیت} & \textbf{مراحل پیاده سازی} & \textbf{تست و پذیرش} \\\hline
\endhead
76 & بانک استعدادهای TA & پروفایل آموزشی و skill/experience/availability را مدل کن؛ search/filter و scoring service بساز؛ صفحه talent-pool را تکمیل کن؛ تست ranking و privacy بنویس. برای این قابلیت، ابتدا acceptance criteria بنویس، سپس migration/schema، service، API، UI، audit، test و documentation را در همان PR انجام بده. & Unit + integration + e2e؛ تست permission، validation، audit و حالت خطا. \\ \hline
77 & ثبت مهارت ها & پروفایل آموزشی و skill/experience/availability را مدل کن؛ search/filter و scoring service بساز؛ صفحه talent-pool را تکمیل کن؛ تست ranking و privacy بنویس. برای این قابلیت، ابتدا acceptance criteria بنویس، سپس migration/schema، service، API، UI، audit، test و documentation را در همان PR انجام بده. & Unit + integration + e2e؛ تست permission، validation، audit و حالت خطا. \\ \hline
78 & ثبت درس های گذرانده شده & پروفایل آموزشی و skill/experience/availability را مدل کن؛ search/filter و scoring service بساز؛ صفحه talent-pool را تکمیل کن؛ تست ranking و privacy بنویس. برای این قابلیت، ابتدا acceptance criteria بنویس، سپس migration/schema، service، API، UI، audit، test و documentation را در همان PR انجام بده. & Unit + integration + e2e؛ تست permission، validation، audit و حالت خطا. \\ \hline
79 & ثبت نمره درس های مهم & پروفایل آموزشی و skill/experience/availability را مدل کن؛ search/filter و scoring service بساز؛ صفحه talent-pool را تکمیل کن؛ تست ranking و privacy بنویس. برای این قابلیت، ابتدا acceptance criteria بنویس، سپس migration/schema، service، API، UI، audit، test و documentation را در همان PR انجام بده. & Unit + integration + e2e؛ تست permission، validation، audit و حالت خطا. \\ \hline
80 & سابقه TA قبلی & پروفایل آموزشی و skill/experience/availability را مدل کن؛ search/filter و scoring service بساز؛ صفحه talent-pool را تکمیل کن؛ تست ranking و privacy بنویس. برای این قابلیت، ابتدا acceptance criteria بنویس، سپس migration/schema، service، API، UI، audit، test و documentation را در همان PR انجام بده. & Unit + integration + e2e؛ تست permission، validation، audit و حالت خطا. \\ \hline
81 & سابقه Head TA & پروفایل آموزشی و skill/experience/availability را مدل کن؛ search/filter و scoring service بساز؛ صفحه talent-pool را تکمیل کن؛ تست ranking و privacy بنویس. برای این قابلیت، ابتدا acceptance criteria بنویس، سپس migration/schema، service، API، UI، audit، test و documentation را در همان PR انجام بده. & Unit + integration + e2e؛ تست permission، validation، audit و حالت خطا. \\ \hline
82 & سابقه تصحیح تمرین & پروفایل آموزشی و skill/experience/availability را مدل کن؛ search/filter و scoring service بساز؛ صفحه talent-pool را تکمیل کن؛ تست ranking و privacy بنویس. برای این قابلیت، ابتدا acceptance criteria بنویس، سپس migration/schema، service، API، UI، audit، test و documentation را در همان PR انجام بده. & Unit + integration + e2e؛ تست permission، validation، audit و حالت خطا. \\ \hline
83 & سابقه برگزاری جلسه حل تمرین & پروفایل آموزشی و skill/experience/availability را مدل کن؛ search/filter و scoring service بساز؛ صفحه talent-pool را تکمیل کن؛ تست ranking و privacy بنویس. برای این قابلیت، ابتدا acceptance criteria بنویس، سپس migration/schema، service، API، UI، audit، test و documentation را در همان PR انجام بده. & Unit + integration + e2e؛ تست permission، validation، audit و حالت خطا. \\ \hline
84 & امتیاز ارزیابی قبلی & پروفایل آموزشی و skill/experience/availability را مدل کن؛ search/filter و scoring service بساز؛ صفحه talent-pool را تکمیل کن؛ تست ranking و privacy بنویس. برای این قابلیت، ابتدا acceptance criteria بنویس، سپس migration/schema، service، API، UI، audit، test و documentation را در همان PR انجام بده. & Unit + integration + e2e؛ تست permission، validation، audit و حالت خطا. \\ \hline
85 & availability هفتگی & پروفایل آموزشی و skill/experience/availability را مدل کن؛ search/filter و scoring service بساز؛ صفحه talent-pool را تکمیل کن؛ تست ranking و privacy بنویس. برای این قابلیت، ابتدا acceptance criteria بنویس، سپس migration/schema، service، API، UI، audit، test و documentation را در همان PR انجام بده. & Unit + integration + e2e؛ تست permission، validation، audit و حالت خطا. \\ \hline
86 & professor favorite list & پروفایل آموزشی و skill/experience/availability را مدل کن؛ search/filter و scoring service بساز؛ صفحه talent-pool را تکمیل کن؛ تست ranking و privacy بنویس. برای این قابلیت، ابتدا acceptance criteria بنویس، سپس migration/schema، service، API، UI، audit، test و documentation را در همان PR انجام بده. & Unit + integration + e2e؛ تست permission، validation، audit و حالت خطا. \\ \hline
87 & جستجوی دانشجو براساس مهارت & پروفایل آموزشی و skill/experience/availability را مدل کن؛ search/filter و scoring service بساز؛ صفحه talent-pool را تکمیل کن؛ تست ranking و privacy بنویس. برای این قابلیت، ابتدا acceptance criteria بنویس، سپس migration/schema، service، API، UI، audit، test و documentation را در همان PR انجام بده. & Unit + integration + e2e؛ تست permission، validation، audit و حالت خطا. \\ \hline
88 & پیشنهاد خودکار بهترین افراد & پروفایل آموزشی و skill/experience/availability را مدل کن؛ search/filter و scoring service بساز؛ صفحه talent-pool را تکمیل کن؛ تست ranking و privacy بنویس. برای این قابلیت، ابتدا acceptance criteria بنویس، سپس migration/schema، service، API، UI، audit، test و documentation را در همان PR انجام بده. & Unit + integration + e2e؛ تست permission، validation، audit و حالت خطا. \\ \hline
89 & مقایسه چند کاندید & پروفایل آموزشی و skill/experience/availability را مدل کن؛ search/filter و scoring service بساز؛ صفحه talent-pool را تکمیل کن؛ تست ranking و privacy بنویس. برای این قابلیت، ابتدا acceptance criteria بنویس، سپس migration/schema، service، API، UI، audit، test و documentation را در همان PR انجام بده. & Unit + integration + e2e؛ تست permission، validation، audit و حالت خطا. \\ \hline
90 & خروجی رزومه TA & پروفایل آموزشی و skill/experience/availability را مدل کن؛ search/filter و scoring service بساز؛ صفحه talent-pool را تکمیل کن؛ تست ranking و privacy بنویس. برای این قابلیت، ابتدا acceptance criteria بنویس، سپس migration/schema، service، API، UI، audit، test و documentation را در همان PR انجام بده. & Unit + integration + e2e؛ تست permission، validation، audit و حالت خطا. \\ \hline
\end{longtable}
\subsection{F) کلاس، جلسات و حضور}

\begin{longtable}{|>{\raggedleft\arraybackslash}p{0.065\textwidth}|>{\raggedleft\arraybackslash}p{0.21\textwidth}|>{\raggedleft\arraybackslash}p{0.42\textwidth}|>{\raggedleft\arraybackslash}p{0.18\textwidth}|}
\hline
\textbf{شماره} & \textbf{قابلیت} & \textbf{مراحل پیاده سازی} & \textbf{تست و پذیرش} \\\hline
\endfirsthead
\hline
\textbf{شماره} & \textbf{قابلیت} & \textbf{مراحل پیاده سازی} & \textbf{تست و پذیرش} \\\hline
\endhead
91 & ساخت جلسه رفع اشکال & مدل session/registration/attendance را کامل کن؛ conflict و capacity را validate کن؛ UI ثبت نام و ICS بساز؛ e2e register/cancel/attendance بنویس. برای این قابلیت، ابتدا acceptance criteria بنویس، سپس migration/schema، service، API، UI، audit، test و documentation را در همان PR انجام بده. & Unit + integration + e2e؛ تست permission، validation، audit و حالت خطا. \\ \hline
92 & لینک Google Meet دستی & مدل session/registration/attendance را کامل کن؛ conflict و capacity را validate کن؛ UI ثبت نام و ICS بساز؛ e2e register/cancel/attendance بنویس. برای این قابلیت، ابتدا acceptance criteria بنویس، سپس migration/schema، service، API، UI، audit، test و documentation را در همان PR انجام بده. & Unit + integration + e2e؛ تست permission، validation، audit و حالت خطا. \\ \hline
93 & مکان فیزیکی جلسه & مدل session/registration/attendance را کامل کن؛ conflict و capacity را validate کن؛ UI ثبت نام و ICS بساز؛ e2e register/cancel/attendance بنویس. برای این قابلیت، ابتدا acceptance criteria بنویس، سپس migration/schema، service، API، UI، audit، test و documentation را در همان PR انجام بده. & Unit + integration + e2e؛ تست permission، validation، audit و حالت خطا. \\ \hline
94 & ظرفیت جلسه & مدل session/registration/attendance را کامل کن؛ conflict و capacity را validate کن؛ UI ثبت نام و ICS بساز؛ e2e register/cancel/attendance بنویس. برای این قابلیت، ابتدا acceptance criteria بنویس، سپس migration/schema، service، API، UI، audit، test و documentation را در همان PR انجام بده. & Unit + integration + e2e؛ تست permission، validation، audit و حالت خطا. \\ \hline
95 & ثبت نام دانشجو در جلسه & مدل session/registration/attendance را کامل کن؛ conflict و capacity را validate کن؛ UI ثبت نام و ICS بساز؛ e2e register/cancel/attendance بنویس. برای این قابلیت، ابتدا acceptance criteria بنویس، سپس migration/schema، service، API، UI، audit، test و documentation را در همان PR انجام بده. & Unit + integration + e2e؛ تست permission، validation، audit و حالت خطا. \\ \hline
96 & cancel registration & مدل session/registration/attendance را کامل کن؛ conflict و capacity را validate کن؛ UI ثبت نام و ICS بساز؛ e2e register/cancel/attendance بنویس. برای این قابلیت، ابتدا acceptance criteria بنویس، سپس migration/schema، service، API، UI، audit، test و documentation را در همان PR انجام بده. & Unit + integration + e2e؛ تست permission، validation، audit و حالت خطا. \\ \hline
97 & waitlist & مدل session/registration/attendance را کامل کن؛ conflict و capacity را validate کن؛ UI ثبت نام و ICS بساز؛ e2e register/cancel/attendance بنویس. برای این قابلیت، ابتدا acceptance criteria بنویس، سپس migration/schema، service، API، UI، audit، test و documentation را در همان PR انجام بده. & Unit + integration + e2e؛ تست permission، validation، audit و حالت خطا. \\ \hline
98 & recurring session & مدل session/registration/attendance را کامل کن؛ conflict و capacity را validate کن؛ UI ثبت نام و ICS بساز؛ e2e register/cancel/attendance بنویس. برای این قابلیت، ابتدا acceptance criteria بنویس، سپس migration/schema، service، API، UI، audit، test و documentation را در همان PR انجام بده. & Unit + integration + e2e؛ تست permission، validation، audit و حالت خطا. \\ \hline
99 & reschedule/cancel session & مدل session/registration/attendance را کامل کن؛ conflict و capacity را validate کن؛ UI ثبت نام و ICS بساز؛ e2e register/cancel/attendance بنویس. برای این قابلیت، ابتدا acceptance criteria بنویس، سپس migration/schema، service، API، UI، audit، test و documentation را در همان PR انجام بده. & Unit + integration + e2e؛ تست permission، validation، audit و حالت خطا. \\ \hline
100 & export ICS برای تقویم & مدل session/registration/attendance را کامل کن؛ conflict و capacity را validate کن؛ UI ثبت نام و ICS بساز؛ e2e register/cancel/attendance بنویس. برای این قابلیت، ابتدا acceptance criteria بنویس، سپس migration/schema، service، API، UI، audit، test و documentation را در همان PR انجام بده. & Unit + integration + e2e؛ تست permission، validation، audit و حالت خطا. \\ \hline
\end{longtable}
\subsection{G) تکلیف، تحویل، تصحیح}

\begin{longtable}{|>{\raggedleft\arraybackslash}p{0.065\textwidth}|>{\raggedleft\arraybackslash}p{0.21\textwidth}|>{\raggedleft\arraybackslash}p{0.42\textwidth}|>{\raggedleft\arraybackslash}p{0.18\textwidth}|}
\hline
\textbf{شماره} & \textbf{قابلیت} & \textbf{مراحل پیاده سازی} & \textbf{تست و پذیرش} \\\hline
\endfirsthead
\hline
\textbf{شماره} & \textbf{قابلیت} & \textbf{مراحل پیاده سازی} & \textbf{تست و پذیرش} \\\hline
\endhead
101 & تعریف تکلیف توسط استاد/TA & مدل Assignment/Submission/Rubric را اضافه کن؛ upload و late policy را enforce کن؛ UI تحویل/تصحیح بساز؛ تست deadline، permission و file access بنویس. برای این قابلیت، ابتدا acceptance criteria بنویس، سپس migration/schema، service، API، UI، audit، test و documentation را در همان PR انجام بده. & Unit + integration + e2e؛ تست permission، validation، audit و حالت خطا. سناریوهای deadline، import/rollback، visibility و score bounds را تست کن. \\ \hline
102 & deadline تکلیف & مدل Assignment/Submission/Rubric را اضافه کن؛ upload و late policy را enforce کن؛ UI تحویل/تصحیح بساز؛ تست deadline، permission و file access بنویس. برای این قابلیت، ابتدا acceptance criteria بنویس، سپس migration/schema، service، API، UI، audit، test و documentation را در همان PR انجام بده. & Unit + integration + e2e؛ تست permission، validation، audit و حالت خطا. سناریوهای deadline، import/rollback، visibility و score bounds را تست کن. \\ \hline
103 & upload فایل تکلیف توسط دانشجو & مدل Assignment/Submission/Rubric را اضافه کن؛ upload و late policy را enforce کن؛ UI تحویل/تصحیح بساز؛ تست deadline، permission و file access بنویس. برای این قابلیت، ابتدا acceptance criteria بنویس، سپس migration/schema، service، API، UI، audit، test و documentation را در همان PR انجام بده. & Unit + integration + e2e؛ تست permission، validation، audit و حالت خطا. سناریوهای deadline، import/rollback، visibility و score bounds را تست کن. \\ \hline
104 & text submission & مدل Assignment/Submission/Rubric را اضافه کن؛ upload و late policy را enforce کن؛ UI تحویل/تصحیح بساز؛ تست deadline، permission و file access بنویس. برای این قابلیت، ابتدا acceptance criteria بنویس، سپس migration/schema، service، API، UI، audit، test و documentation را در همان PR انجام بده. & Unit + integration + e2e؛ تست permission، validation، audit و حالت خطا. سناریوهای deadline، import/rollback، visibility و score bounds را تست کن. \\ \hline
105 & multiple file submission & مدل Assignment/Submission/Rubric را اضافه کن؛ upload و late policy را enforce کن؛ UI تحویل/تصحیح بساز؛ تست deadline، permission و file access بنویس. برای این قابلیت، ابتدا acceptance criteria بنویس، سپس migration/schema، service، API، UI، audit، test و documentation را در همان PR انجام بده. & Unit + integration + e2e؛ تست permission، validation، audit و حالت خطا. سناریوهای deadline، import/rollback، visibility و score bounds را تست کن. \\ \hline
106 & late submission policy & مدل Assignment/Submission/Rubric را اضافه کن؛ upload و late policy را enforce کن؛ UI تحویل/تصحیح بساز؛ تست deadline، permission و file access بنویس. برای این قابلیت، ابتدا acceptance criteria بنویس، سپس migration/schema، service، API، UI، audit، test و documentation را در همان PR انجام بده. & Unit + integration + e2e؛ تست permission، validation، audit و حالت خطا. سناریوهای deadline، import/rollback، visibility و score bounds را تست کن. \\ \hline
107 & grace period & مدل Assignment/Submission/Rubric را اضافه کن؛ upload و late policy را enforce کن؛ UI تحویل/تصحیح بساز؛ تست deadline، permission و file access بنویس. برای این قابلیت، ابتدا acceptance criteria بنویس، سپس migration/schema، service، API، UI، audit، test و documentation را در همان PR انجام بده. & Unit + integration + e2e؛ تست permission، validation، audit و حالت خطا. سناریوهای deadline، import/rollback، visibility و score bounds را تست کن. \\ \hline
108 & گروه بندی تکلیف تیمی & مدل Assignment/Submission/Rubric را اضافه کن؛ upload و late policy را enforce کن؛ UI تحویل/تصحیح بساز؛ تست deadline، permission و file access بنویس. برای این قابلیت، ابتدا acceptance criteria بنویس، سپس migration/schema، service، API، UI، audit، test و documentation را در همان PR انجام بده. & Unit + integration + e2e؛ تست permission، validation، audit و حالت خطا. سناریوهای deadline، import/rollback، visibility و score bounds را تست کن. \\ \hline
109 & assign تصحیح به TA & مدل Assignment/Submission/Rubric را اضافه کن؛ upload و late policy را enforce کن؛ UI تحویل/تصحیح بساز؛ تست deadline، permission و file access بنویس. برای این قابلیت، ابتدا acceptance criteria بنویس، سپس migration/schema، service، API، UI، audit، test و documentation را در همان PR انجام بده. & Unit + integration + e2e؛ تست permission، validation، audit و حالت خطا. سناریوهای deadline، import/rollback، visibility و score bounds را تست کن. \\ \hline
110 & rubric تصحیح & مدل Assignment/Submission/Rubric را اضافه کن؛ upload و late policy را enforce کن؛ UI تحویل/تصحیح بساز؛ تست deadline، permission و file access بنویس. برای این قابلیت، ابتدا acceptance criteria بنویس، سپس migration/schema، service، API، UI، audit، test و documentation را در همان PR انجام بده. & Unit + integration + e2e؛ تست permission، validation، audit و حالت خطا. سناریوهای deadline، import/rollback، visibility و score bounds را تست کن. \\ \hline
111 & comment روی submission & مدل Assignment/Submission/Rubric را اضافه کن؛ upload و late policy را enforce کن؛ UI تحویل/تصحیح بساز؛ تست deadline، permission و file access بنویس. برای این قابلیت، ابتدا acceptance criteria بنویس، سپس migration/schema، service، API، UI، audit، test و documentation را در همان PR انجام بده. & Unit + integration + e2e؛ تست permission، validation، audit و حالت خطا. سناریوهای deadline، import/rollback، visibility و score bounds را تست کن. \\ \hline
112 & فایل feedback & مدل Assignment/Submission/Rubric را اضافه کن؛ upload و late policy را enforce کن؛ UI تحویل/تصحیح بساز؛ تست deadline، permission و file access بنویس. برای این قابلیت، ابتدا acceptance criteria بنویس، سپس migration/schema، service، API، UI، audit، test و documentation را در همان PR انجام بده. & Unit + integration + e2e؛ تست permission، validation، audit و حالت خطا. سناریوهای deadline، import/rollback، visibility و score bounds را تست کن. \\ \hline
113 & نمره دهی مرحله ای & مدل Assignment/Submission/Rubric را اضافه کن؛ upload و late policy را enforce کن؛ UI تحویل/تصحیح بساز؛ تست deadline، permission و file access بنویس. برای این قابلیت، ابتدا acceptance criteria بنویس، سپس migration/schema، service، API، UI، audit، test و documentation را در همان PR انجام بده. & Unit + integration + e2e؛ تست permission، validation، audit و حالت خطا. سناریوهای deadline، import/rollback، visibility و score bounds را تست کن. \\ \hline
114 & publish feedback & مدل Assignment/Submission/Rubric را اضافه کن؛ upload و late policy را enforce کن؛ UI تحویل/تصحیح بساز؛ تست deadline، permission و file access بنویس. برای این قابلیت، ابتدا acceptance criteria بنویس، سپس migration/schema، service، API، UI، audit، test و documentation را در همان PR انجام بده. & Unit + integration + e2e؛ تست permission، validation، audit و حالت خطا. سناریوهای deadline، import/rollback، visibility و score bounds را تست کن. \\ \hline
115 & plagiarism/similarity report & مدل Assignment/Submission/Rubric را اضافه کن؛ upload و late policy را enforce کن؛ UI تحویل/تصحیح بساز؛ تست deadline، permission و file access بنویس. برای این قابلیت، ابتدا acceptance criteria بنویس، سپس migration/schema، service، API، UI، audit، test و documentation را در همان PR انجام بده. & Unit + integration + e2e؛ تست permission، validation، audit و حالت خطا. سناریوهای deadline، import/rollback، visibility و score bounds را تست کن. \\ \hline
116 & درخواست تمدید مهلت & مدل Assignment/Submission/Rubric را اضافه کن؛ upload و late policy را enforce کن؛ UI تحویل/تصحیح بساز؛ تست deadline، permission و file access بنویس. برای این قابلیت، ابتدا acceptance criteria بنویس، سپس migration/schema، service، API، UI، audit، test و documentation را در همان PR انجام بده. & Unit + integration + e2e؛ تست permission، validation، audit و حالت خطا. سناریوهای deadline، import/rollback، visibility و score bounds را تست کن. \\ \hline
117 & approval برای extension & مدل Assignment/Submission/Rubric را اضافه کن؛ upload و late policy را enforce کن؛ UI تحویل/تصحیح بساز؛ تست deadline، permission و file access بنویس. برای این قابلیت، ابتدا acceptance criteria بنویس، سپس migration/schema، service، API، UI، audit، test و documentation را در همان PR انجام بده. & Unit + integration + e2e؛ تست permission، validation، audit و حالت خطا. سناریوهای deadline، import/rollback، visibility و score bounds را تست کن. \\ \hline
118 & export submissions & مدل Assignment/Submission/Rubric را اضافه کن؛ upload و late policy را enforce کن؛ UI تحویل/تصحیح بساز؛ تست deadline، permission و file access بنویس. برای این قابلیت، ابتدا acceptance criteria بنویس، سپس migration/schema، service، API، UI، audit، test و documentation را در همان PR انجام بده. & Unit + integration + e2e؛ تست permission، validation، audit و حالت خطا. سناریوهای deadline، import/rollback، visibility و score bounds را تست کن. \\ \hline
119 & bulk download submissions & مدل Assignment/Submission/Rubric را اضافه کن؛ upload و late policy را enforce کن؛ UI تحویل/تصحیح بساز؛ تست deadline، permission و file access بنویس. برای این قابلیت، ابتدا acceptance criteria بنویس، سپس migration/schema، service، API، UI، audit، test و documentation را در همان PR انجام بده. & Unit + integration + e2e؛ تست permission، validation، audit و حالت خطا. سناریوهای deadline، import/rollback، visibility و score bounds را تست کن. \\ \hline
120 & submission history & مدل Assignment/Submission/Rubric را اضافه کن؛ upload و late policy را enforce کن؛ UI تحویل/تصحیح بساز؛ تست deadline، permission و file access بنویس. برای این قابلیت، ابتدا acceptance criteria بنویس، سپس migration/schema، service، API، UI، audit، test و documentation را در همان PR انجام بده. & Unit + integration + e2e؛ تست permission، validation، audit و حالت خطا. سناریوهای deadline، import/rollback، visibility و score bounds را تست کن. \\ \hline
\end{longtable}
\subsection{H) Gradebook}

\begin{longtable}{|>{\raggedleft\arraybackslash}p{0.065\textwidth}|>{\raggedleft\arraybackslash}p{0.21\textwidth}|>{\raggedleft\arraybackslash}p{0.42\textwidth}|>{\raggedleft\arraybackslash}p{0.18\textwidth}|}
\hline
\textbf{شماره} & \textbf{قابلیت} & \textbf{مراحل پیاده سازی} & \textbf{تست و پذیرش} \\\hline
\endfirsthead
\hline
\textbf{شماره} & \textbf{قابلیت} & \textbf{مراحل پیاده سازی} & \textbf{تست و پذیرش} \\\hline
\endhead
121 & دسته بندی نمرات & GradeCategory/Item/Record را با constraint و transaction کامل کن؛ import/export و publish/lock را بساز؛ UI gradebook را role-based کن؛ تست score bounds، visibility و rollback بنویس. برای این قابلیت، ابتدا acceptance criteria بنویس، سپس migration/schema، service، API، UI، audit، test و documentation را در همان PR انجام بده. & Unit + integration + e2e؛ تست permission، validation، audit و حالت خطا. سناریوهای deadline، import/rollback، visibility و score bounds را تست کن. \\ \hline
122 & وزن دهی categories & GradeCategory/Item/Record را با constraint و transaction کامل کن؛ import/export و publish/lock را بساز؛ UI gradebook را role-based کن؛ تست score bounds، visibility و rollback بنویس. برای این قابلیت، ابتدا acceptance criteria بنویس، سپس migration/schema، service، API، UI، audit، test و documentation را در همان PR انجام بده. & Unit + integration + e2e؛ تست permission، validation، audit و حالت خطا. سناریوهای deadline، import/rollback، visibility و score bounds را تست کن. \\ \hline
123 & GradeItem & GradeCategory/Item/Record را با constraint و transaction کامل کن؛ import/export و publish/lock را بساز؛ UI gradebook را role-based کن؛ تست score bounds، visibility و rollback بنویس. برای این قابلیت، ابتدا acceptance criteria بنویس، سپس migration/schema، service، API، UI، audit، test و documentation را در همان PR انجام بده. & Unit + integration + e2e؛ تست permission، validation، audit و حالت خطا. سناریوهای deadline، import/rollback، visibility و score bounds را تست کن. \\ \hline
124 & GradeRecord & GradeCategory/Item/Record را با constraint و transaction کامل کن؛ import/export و publish/lock را بساز؛ UI gradebook را role-based کن؛ تست score bounds، visibility و rollback بنویس. برای این قابلیت، ابتدا acceptance criteria بنویس، سپس migration/schema، service، API، UI، audit، test و documentation را در همان PR انجام بده. & Unit + integration + e2e؛ تست permission، validation، audit و حالت خطا. سناریوهای deadline، import/rollback، visibility و score bounds را تست کن. \\ \hline
125 & edit grade & GradeCategory/Item/Record را با constraint و transaction کامل کن؛ import/export و publish/lock را بساز؛ UI gradebook را role-based کن؛ تست score bounds، visibility و rollback بنویس. برای این قابلیت، ابتدا acceptance criteria بنویس، سپس migration/schema، service، API، UI، audit، test و documentation را در همان PR انجام بده. & Unit + integration + e2e؛ تست permission، validation، audit و حالت خطا. سناریوهای deadline، import/rollback، visibility و score bounds را تست کن. \\ \hline
126 & grade change history & GradeCategory/Item/Record را با constraint و transaction کامل کن؛ import/export و publish/lock را بساز؛ UI gradebook را role-based کن؛ تست score bounds، visibility و rollback بنویس. برای این قابلیت، ابتدا acceptance criteria بنویس، سپس migration/schema، service، API، UI، audit، test و documentation را در همان PR انجام بده. & Unit + integration + e2e؛ تست permission، validation، audit و حالت خطا. سناریوهای deadline، import/rollback، visibility و score bounds را تست کن. \\ \hline
127 & lock grade & GradeCategory/Item/Record را با constraint و transaction کامل کن؛ import/export و publish/lock را بساز؛ UI gradebook را role-based کن؛ تست score bounds، visibility و rollback بنویس. برای این قابلیت، ابتدا acceptance criteria بنویس، سپس migration/schema، service، API، UI، audit، test و documentation را در همان PR انجام بده. & Unit + integration + e2e؛ تست permission، validation، audit و حالت خطا. سناریوهای deadline، import/rollback، visibility و score bounds را تست کن. \\ \hline
128 & publish grade & GradeCategory/Item/Record را با constraint و transaction کامل کن؛ import/export و publish/lock را بساز؛ UI gradebook را role-based کن؛ تست score bounds، visibility و rollback بنویس. برای این قابلیت، ابتدا acceptance criteria بنویس، سپس migration/schema، service، API، UI، audit، test و documentation را در همان PR انجام بده. & Unit + integration + e2e؛ تست permission، validation، audit و حالت خطا. سناریوهای deadline، import/rollback، visibility و score bounds را تست کن. \\ \hline
129 & unpublish grade & GradeCategory/Item/Record را با constraint و transaction کامل کن؛ import/export و publish/lock را بساز؛ UI gradebook را role-based کن؛ تست score bounds، visibility و rollback بنویس. برای این قابلیت، ابتدا acceptance criteria بنویس، سپس migration/schema، service، API، UI، audit، test و documentation را در همان PR انجام بده. & Unit + integration + e2e؛ تست permission، validation، audit و حالت خطا. سناریوهای deadline، import/rollback، visibility و score bounds را تست کن. \\ \hline
130 & import نمره از Excel & GradeCategory/Item/Record را با constraint و transaction کامل کن؛ import/export و publish/lock را بساز؛ UI gradebook را role-based کن؛ تست score bounds، visibility و rollback بنویس. برای این قابلیت، ابتدا acceptance criteria بنویس، سپس migration/schema، service، API، UI، audit، test و documentation را در همان PR انجام بده. & Unit + integration + e2e؛ تست permission، validation، audit و حالت خطا. سناریوهای deadline، import/rollback، visibility و score bounds را تست کن. \\ \hline
131 & preview import & GradeCategory/Item/Record را با constraint و transaction کامل کن؛ import/export و publish/lock را بساز؛ UI gradebook را role-based کن؛ تست score bounds، visibility و rollback بنویس. برای این قابلیت، ابتدا acceptance criteria بنویس، سپس migration/schema، service، API، UI، audit، test و documentation را در همان PR انجام بده. & Unit + integration + e2e؛ تست permission، validation، audit و حالت خطا. سناریوهای deadline، import/rollback، visibility و score bounds را تست کن. \\ \hline
132 & rollback import & GradeCategory/Item/Record را با constraint و transaction کامل کن؛ import/export و publish/lock را بساز؛ UI gradebook را role-based کن؛ تست score bounds، visibility و rollback بنویس. برای این قابلیت، ابتدا acceptance criteria بنویس، سپس migration/schema، service، API، UI، audit، test و documentation را در همان PR انجام بده. & Unit + integration + e2e؛ تست permission، validation، audit و حالت خطا. سناریوهای deadline، import/rollback، visibility و score bounds را تست کن. \\ \hline
133 & validation خطاهای import & GradeCategory/Item/Record را با constraint و transaction کامل کن؛ import/export و publish/lock را بساز؛ UI gradebook را role-based کن؛ تست score bounds، visibility و rollback بنویس. برای این قابلیت، ابتدا acceptance criteria بنویس، سپس migration/schema، service، API، UI، audit، test و documentation را در همان PR انجام بده. & Unit + integration + e2e؛ تست permission، validation، audit و حالت خطا. سناریوهای deadline، import/rollback، visibility و score bounds را تست کن. \\ \hline
134 & export gradebook CSV & GradeCategory/Item/Record را با constraint و transaction کامل کن؛ import/export و publish/lock را بساز؛ UI gradebook را role-based کن؛ تست score bounds، visibility و rollback بنویس. برای این قابلیت، ابتدا acceptance criteria بنویس، سپس migration/schema، service، API، UI، audit، test و documentation را در همان PR انجام بده. & Unit + integration + e2e؛ تست permission، validation، audit و حالت خطا. سناریوهای deadline، import/rollback، visibility و score bounds را تست کن. \\ \hline
135 & export gradebook XLSX & GradeCategory/Item/Record را با constraint و transaction کامل کن؛ import/export و publish/lock را بساز؛ UI gradebook را role-based کن؛ تست score bounds، visibility و rollback بنویس. برای این قابلیت، ابتدا acceptance criteria بنویس، سپس migration/schema، service، API، UI، audit، test و documentation را در همان PR انجام بده. & Unit + integration + e2e؛ تست permission، validation، audit و حالت خطا. سناریوهای deadline، import/rollback، visibility و score bounds را تست کن. \\ \hline
136 & محاسبه نمره نهایی & GradeCategory/Item/Record را با constraint و transaction کامل کن؛ import/export و publish/lock را بساز؛ UI gradebook را role-based کن؛ تست score bounds، visibility و rollback بنویس. برای این قابلیت، ابتدا acceptance criteria بنویس، سپس migration/schema، service، API، UI، audit، test و documentation را در همان PR انجام بده. & Unit + integration + e2e؛ تست permission، validation، audit و حالت خطا. سناریوهای deadline، import/rollback، visibility و score bounds را تست کن. \\ \hline
137 & curve/bonus policy & GradeCategory/Item/Record را با constraint و transaction کامل کن؛ import/export و publish/lock را بساز؛ UI gradebook را role-based کن؛ تست score bounds، visibility و rollback بنویس. برای این قابلیت، ابتدا acceptance criteria بنویس، سپس migration/schema، service، API، UI، audit، test و documentation را در همان PR انجام بده. & Unit + integration + e2e؛ تست permission، validation، audit و حالت خطا. سناریوهای deadline، import/rollback، visibility و score bounds را تست کن. \\ \hline
138 & grade appeal & GradeCategory/Item/Record را با constraint و transaction کامل کن؛ import/export و publish/lock را بساز؛ UI gradebook را role-based کن؛ تست score bounds، visibility و rollback بنویس. برای این قابلیت، ابتدا acceptance criteria بنویس، سپس migration/schema، service، API، UI، audit، test و documentation را در همان PR انجام بده. & Unit + integration + e2e؛ تست permission، validation، audit و حالت خطا. سناریوهای deadline، import/rollback، visibility و score bounds را تست کن. \\ \hline
139 & پاسخ به اعتراض نمره & GradeCategory/Item/Record را با constraint و transaction کامل کن؛ import/export و publish/lock را بساز؛ UI gradebook را role-based کن؛ تست score bounds، visibility و rollback بنویس. برای این قابلیت، ابتدا acceptance criteria بنویس، سپس migration/schema، service، API، UI، audit، test و documentation را در همان PR انجام بده. & Unit + integration + e2e؛ تست permission، validation، audit و حالت خطا. سناریوهای deadline، import/rollback، visibility و score bounds را تست کن. \\ \hline
140 & گزارش عملکرد تصحیح TA & GradeCategory/Item/Record را با constraint و transaction کامل کن؛ import/export و publish/lock را بساز؛ UI gradebook را role-based کن؛ تست score bounds، visibility و rollback بنویس. برای این قابلیت، ابتدا acceptance criteria بنویس، سپس migration/schema، service، API، UI، audit، test و documentation را در همان PR انجام بده. & Unit + integration + e2e؛ تست permission، validation، audit و حالت خطا. سناریوهای deadline، import/rollback، visibility و score bounds را تست کن. \\ \hline
\end{longtable}
\subsection{I) پیام رسانی}

\begin{longtable}{|>{\raggedleft\arraybackslash}p{0.065\textwidth}|>{\raggedleft\arraybackslash}p{0.21\textwidth}|>{\raggedleft\arraybackslash}p{0.42\textwidth}|>{\raggedleft\arraybackslash}p{0.18\textwidth}|}
\hline
\textbf{شماره} & \textbf{قابلیت} & \textbf{مراحل پیاده سازی} & \textbf{تست و پذیرش} \\\hline
\endfirsthead
\hline
\textbf{شماره} & \textbf{قابلیت} & \textbf{مراحل پیاده سازی} & \textbf{تست و پذیرش} \\\hline
\endhead
141 & thread عمومی درس & Thread/Message/Participant را با permission و sanitize کامل کن؛ unread/attachment/escalation اضافه کن؛ UI پیام را بهبود بده؛ تست XSS و unauthorized access بنویس. برای این قابلیت، ابتدا acceptance criteria بنویس، سپس migration/schema، service، API، UI، audit، test و documentation را در همان PR انجام بده. & Unit + integration + e2e؛ تست permission، validation، audit و حالت خطا. \\ \hline
142 & thread خصوصی دانشجو با TA & Thread/Message/Participant را با permission و sanitize کامل کن؛ unread/attachment/escalation اضافه کن؛ UI پیام را بهبود بده؛ تست XSS و unauthorized access بنویس. برای این قابلیت، ابتدا acceptance criteria بنویس، سپس migration/schema، service، API، UI، audit، test و documentation را در همان PR انجام بده. & Unit + integration + e2e؛ تست permission، validation، audit و حالت خطا. \\ \hline
143 & thread دانشجو با استاد & Thread/Message/Participant را با permission و sanitize کامل کن؛ unread/attachment/escalation اضافه کن؛ UI پیام را بهبود بده؛ تست XSS و unauthorized access بنویس. برای این قابلیت، ابتدا acceptance criteria بنویس، سپس migration/schema، service، API، UI، audit، test و documentation را در همان PR انجام بده. & Unit + integration + e2e؛ تست permission، validation، audit و حالت خطا. \\ \hline
144 & thread staff-only & Thread/Message/Participant را با permission و sanitize کامل کن؛ unread/attachment/escalation اضافه کن؛ UI پیام را بهبود بده؛ تست XSS و unauthorized access بنویس. برای این قابلیت، ابتدا acceptance criteria بنویس، سپس migration/schema، service، API، UI، audit، test و documentation را در همان PR انجام بده. & Unit + integration + e2e؛ تست permission، validation، audit و حالت خطا. \\ \hline
145 & category پیام & Thread/Message/Participant را با permission و sanitize کامل کن؛ unread/attachment/escalation اضافه کن؛ UI پیام را بهبود بده؛ تست XSS و unauthorized access بنویس. برای این قابلیت، ابتدا acceptance criteria بنویس، سپس migration/schema، service، API، UI، audit، test و documentation را در همان PR انجام بده. & Unit + integration + e2e؛ تست permission، validation، audit و حالت خطا. \\ \hline
146 & priority پیام & Thread/Message/Participant را با permission و sanitize کامل کن؛ unread/attachment/escalation اضافه کن؛ UI پیام را بهبود بده؛ تست XSS و unauthorized access بنویس. برای این قابلیت، ابتدا acceptance criteria بنویس، سپس migration/schema، service، API، UI، audit، test و documentation را در همان PR انجام بده. & Unit + integration + e2e؛ تست permission، validation، audit و حالت خطا. \\ \hline
147 & unread count & Thread/Message/Participant را با permission و sanitize کامل کن؛ unread/attachment/escalation اضافه کن؛ UI پیام را بهبود بده؛ تست XSS و unauthorized access بنویس. برای این قابلیت، ابتدا acceptance criteria بنویس، سپس migration/schema، service، API، UI، audit، test و documentation را در همان PR انجام بده. & Unit + integration + e2e؛ تست permission، validation، audit و حالت خطا. \\ \hline
148 & read receipt & Thread/Message/Participant را با permission و sanitize کامل کن؛ unread/attachment/escalation اضافه کن؛ UI پیام را بهبود بده؛ تست XSS و unauthorized access بنویس. برای این قابلیت، ابتدا acceptance criteria بنویس، سپس migration/schema، service، API، UI، audit، test و documentation را در همان PR انجام بده. & Unit + integration + e2e؛ تست permission، validation، audit و حالت خطا. \\ \hline
149 & attachment & Thread/Message/Participant را با permission و sanitize کامل کن؛ unread/attachment/escalation اضافه کن؛ UI پیام را بهبود بده؛ تست XSS و unauthorized access بنویس. برای این قابلیت، ابتدا acceptance criteria بنویس، سپس migration/schema، service، API، UI، audit، test و documentation را در همان PR انجام بده. & Unit + integration + e2e؛ تست permission، validation، audit و حالت خطا. \\ \hline
150 & sanitize HTML & Thread/Message/Participant را با permission و sanitize کامل کن؛ unread/attachment/escalation اضافه کن؛ UI پیام را بهبود بده؛ تست XSS و unauthorized access بنویس. برای این قابلیت، ابتدا acceptance criteria بنویس، سپس migration/schema، service، API، UI، audit، test و documentation را در همان PR انجام بده. & Unit + integration + e2e؛ تست permission، validation، audit و حالت خطا. \\ \hline
151 & report message & Thread/Message/Participant را با permission و sanitize کامل کن؛ unread/attachment/escalation اضافه کن؛ UI پیام را بهبود بده؛ تست XSS و unauthorized access بنویس. برای این قابلیت، ابتدا acceptance criteria بنویس، سپس migration/schema، service، API، UI، audit، test و documentation را در همان PR انجام بده. & Unit + integration + e2e؛ تست permission، validation، audit و حالت خطا. \\ \hline
152 & mute thread & Thread/Message/Participant را با permission و sanitize کامل کن؛ unread/attachment/escalation اضافه کن؛ UI پیام را بهبود بده؛ تست XSS و unauthorized access بنویس. برای این قابلیت، ابتدا acceptance criteria بنویس، سپس migration/schema، service، API، UI، audit، test و documentation را در همان PR انجام بده. & Unit + integration + e2e؛ تست permission، validation، audit و حالت خطا. \\ \hline
153 & close/reopen thread & Thread/Message/Participant را با permission و sanitize کامل کن؛ unread/attachment/escalation اضافه کن؛ UI پیام را بهبود بده؛ تست XSS و unauthorized access بنویس. برای این قابلیت، ابتدا acceptance criteria بنویس، سپس migration/schema، service، API، UI، audit، test و documentation را در همان PR انجام بده. & Unit + integration + e2e؛ تست permission، validation، audit و حالت خطا. \\ \hline
154 & assign thread به TA & Thread/Message/Participant را با permission و sanitize کامل کن؛ unread/attachment/escalation اضافه کن؛ UI پیام را بهبود بده؛ تست XSS و unauthorized access بنویس. برای این قابلیت، ابتدا acceptance criteria بنویس، سپس migration/schema، service، API، UI، audit، test و documentation را در همان PR انجام بده. & Unit + integration + e2e؛ تست permission، validation، audit و حالت خطا. \\ \hline
155 & escalation به Head TA & Thread/Message/Participant را با permission و sanitize کامل کن؛ unread/attachment/escalation اضافه کن؛ UI پیام را بهبود بده؛ تست XSS و unauthorized access بنویس. برای این قابلیت، ابتدا acceptance criteria بنویس، سپس migration/schema، service، API، UI، audit، test و documentation را در همان PR انجام بده. & Unit + integration + e2e؛ تست permission، validation، audit و حالت خطا. \\ \hline
156 & escalation به professor & Thread/Message/Participant را با permission و sanitize کامل کن؛ unread/attachment/escalation اضافه کن؛ UI پیام را بهبود بده؛ تست XSS و unauthorized access بنویس. برای این قابلیت، ابتدا acceptance criteria بنویس، سپس migration/schema، service، API، UI، audit، test و documentation را در همان PR انجام بده. & Unit + integration + e2e؛ تست permission، validation، audit و حالت خطا. \\ \hline
157 & canned response & Thread/Message/Participant را با permission و sanitize کامل کن؛ unread/attachment/escalation اضافه کن؛ UI پیام را بهبود بده؛ تست XSS و unauthorized access بنویس. برای این قابلیت، ابتدا acceptance criteria بنویس، سپس migration/schema، service، API، UI، audit، test و documentation را در همان PR انجام بده. & Unit + integration + e2e؛ تست permission، validation، audit و حالت خطا. \\ \hline
158 & search messages & Thread/Message/Participant را با permission و sanitize کامل کن؛ unread/attachment/escalation اضافه کن؛ UI پیام را بهبود بده؛ تست XSS و unauthorized access بنویس. برای این قابلیت، ابتدا acceptance criteria بنویس، سپس migration/schema، service، API، UI، audit، test و documentation را در همان PR انجام بده. & Unit + integration + e2e؛ تست permission، validation، audit و حالت خطا. \\ \hline
159 & message audit & Thread/Message/Participant را با permission و sanitize کامل کن؛ unread/attachment/escalation اضافه کن؛ UI پیام را بهبود بده؛ تست XSS و unauthorized access بنویس. برای این قابلیت، ابتدا acceptance criteria بنویس، سپس migration/schema، service، API، UI، audit، test و documentation را در همان PR انجام بده. & Unit + integration + e2e؛ تست permission، validation، audit و حالت خطا. \\ \hline
160 & rate limit پیام & Thread/Message/Participant را با permission و sanitize کامل کن؛ unread/attachment/escalation اضافه کن؛ UI پیام را بهبود بده؛ تست XSS و unauthorized access بنویس. برای این قابلیت، ابتدا acceptance criteria بنویس، سپس migration/schema، service، API، UI، audit، test و documentation را در همان PR انجام بده. & Unit + integration + e2e؛ تست permission، validation، audit و حالت خطا. \\ \hline
\end{longtable}
\subsection{J) نظرسنجی و ارزیابی}

\begin{longtable}{|>{\raggedleft\arraybackslash}p{0.065\textwidth}|>{\raggedleft\arraybackslash}p{0.21\textwidth}|>{\raggedleft\arraybackslash}p{0.42\textwidth}|>{\raggedleft\arraybackslash}p{0.18\textwidth}|}
\hline
\textbf{شماره} & \textbf{قابلیت} & \textbf{مراحل پیاده سازی} & \textbf{تست و پذیرش} \\\hline
\endfirsthead
\hline
\textbf{شماره} & \textbf{قابلیت} & \textbf{مراحل پیاده سازی} & \textbf{تست و پذیرش} \\\hline
\endhead
161 & ساخت survey & Survey/Question/Response/Answer را با anonymity و min response count کامل کن؛ dashboard results بساز؛ export امن اضافه کن؛ تست duplicate و privacy بنویس. برای این قابلیت، ابتدا acceptance criteria بنویس، سپس migration/schema، service، API، UI، audit، test و documentation را در همان PR انجام بده. & Unit + integration + e2e؛ تست permission، validation، audit و حالت خطا. duplicate submission/vote و privacy/anonymity را تست کن. \\ \hline
162 & question builder & Survey/Question/Response/Answer را با anonymity و min response count کامل کن؛ dashboard results بساز؛ export امن اضافه کن؛ تست duplicate و privacy بنویس. برای این قابلیت، ابتدا acceptance criteria بنویس، سپس migration/schema، service، API، UI، audit، test و documentation را در همان PR انجام بده. & Unit + integration + e2e؛ تست permission، validation، audit و حالت خطا. duplicate submission/vote و privacy/anonymity را تست کن. \\ \hline
163 & rating question & Survey/Question/Response/Answer را با anonymity و min response count کامل کن؛ dashboard results بساز؛ export امن اضافه کن؛ تست duplicate و privacy بنویس. برای این قابلیت، ابتدا acceptance criteria بنویس، سپس migration/schema، service، API، UI، audit، test و documentation را در همان PR انجام بده. & Unit + integration + e2e؛ تست permission، validation، audit و حالت خطا. duplicate submission/vote و privacy/anonymity را تست کن. \\ \hline
164 & text question & Survey/Question/Response/Answer را با anonymity و min response count کامل کن؛ dashboard results بساز؛ export امن اضافه کن؛ تست duplicate و privacy بنویس. برای این قابلیت، ابتدا acceptance criteria بنویس، سپس migration/schema، service، API، UI، audit، test و documentation را در همان PR انجام بده. & Unit + integration + e2e؛ تست permission، validation، audit و حالت خطا. duplicate submission/vote و privacy/anonymity را تست کن. \\ \hline
165 & single choice & Survey/Question/Response/Answer را با anonymity و min response count کامل کن؛ dashboard results بساز؛ export امن اضافه کن؛ تست duplicate و privacy بنویس. برای این قابلیت، ابتدا acceptance criteria بنویس، سپس migration/schema، service، API، UI، audit، test و documentation را در همان PR انجام بده. & Unit + integration + e2e؛ تست permission، validation، audit و حالت خطا. duplicate submission/vote و privacy/anonymity را تست کن. \\ \hline
166 & multiple choice & Survey/Question/Response/Answer را با anonymity و min response count کامل کن؛ dashboard results بساز؛ export امن اضافه کن؛ تست duplicate و privacy بنویس. برای این قابلیت، ابتدا acceptance criteria بنویس، سپس migration/schema، service، API، UI، audit، test و documentation را در همان PR انجام بده. & Unit + integration + e2e؛ تست permission، validation، audit و حالت خطا. duplicate submission/vote و privacy/anonymity را تست کن. \\ \hline
167 & boolean question & Survey/Question/Response/Answer را با anonymity و min response count کامل کن؛ dashboard results بساز؛ export امن اضافه کن؛ تست duplicate و privacy بنویس. برای این قابلیت، ابتدا acceptance criteria بنویس، سپس migration/schema، service، API، UI، audit، test و documentation را در همان PR انجام بده. & Unit + integration + e2e؛ تست permission، validation، audit و حالت خطا. duplicate submission/vote و privacy/anonymity را تست کن. \\ \hline
168 & anonymous response & Survey/Question/Response/Answer را با anonymity و min response count کامل کن؛ dashboard results بساز؛ export امن اضافه کن؛ تست duplicate و privacy بنویس. برای این قابلیت، ابتدا acceptance criteria بنویس، سپس migration/schema، service، API، UI، audit، test و documentation را در همان PR انجام بده. & Unit + integration + e2e؛ تست permission، validation، audit و حالت خطا. duplicate submission/vote و privacy/anonymity را تست کن. \\ \hline
169 & respondent hash & Survey/Question/Response/Answer را با anonymity و min response count کامل کن؛ dashboard results بساز؛ export امن اضافه کن؛ تست duplicate و privacy بنویس. برای این قابلیت، ابتدا acceptance criteria بنویس، سپس migration/schema، service، API، UI، audit، test و documentation را در همان PR انجام بده. & Unit + integration + e2e؛ تست permission، validation، audit و حالت خطا. duplicate submission/vote و privacy/anonymity را تست کن. \\ \hline
170 & minimum response count & Survey/Question/Response/Answer را با anonymity و min response count کامل کن؛ dashboard results بساز؛ export امن اضافه کن؛ تست duplicate و privacy بنویس. برای این قابلیت، ابتدا acceptance criteria بنویس، سپس migration/schema، service، API، UI، audit، test و documentation را در همان PR انجام بده. & Unit + integration + e2e؛ تست permission، validation، audit و حالت خطا. duplicate submission/vote و privacy/anonymity را تست کن. \\ \hline
171 & survey result dashboard & Survey/Question/Response/Answer را با anonymity و min response count کامل کن؛ dashboard results بساز؛ export امن اضافه کن؛ تست duplicate و privacy بنویس. برای این قابلیت، ابتدا acceptance criteria بنویس، سپس migration/schema، service، API، UI، audit، test و documentation را در همان PR انجام بده. & Unit + integration + e2e؛ تست permission، validation، audit و حالت خطا. duplicate submission/vote و privacy/anonymity را تست کن. \\ \hline
172 & export survey results & Survey/Question/Response/Answer را با anonymity و min response count کامل کن؛ dashboard results بساز؛ export امن اضافه کن؛ تست duplicate و privacy بنویس. برای این قابلیت، ابتدا acceptance criteria بنویس، سپس migration/schema، service، API، UI، audit، test و documentation را در همان PR انجام بده. & Unit + integration + e2e؛ تست permission، validation، audit و حالت خطا. duplicate submission/vote و privacy/anonymity را تست کن. \\ \hline
173 & TA midterm evaluation & Survey/Question/Response/Answer را با anonymity و min response count کامل کن؛ dashboard results بساز؛ export امن اضافه کن؛ تست duplicate و privacy بنویس. برای این قابلیت، ابتدا acceptance criteria بنویس، سپس migration/schema، service، API، UI، audit، test و documentation را در همان PR انجام بده. & Unit + integration + e2e؛ تست permission، validation، audit و حالت خطا. duplicate submission/vote و privacy/anonymity را تست کن. \\ \hline
174 & TA final evaluation & Survey/Question/Response/Answer را با anonymity و min response count کامل کن؛ dashboard results بساز؛ export امن اضافه کن؛ تست duplicate و privacy بنویس. برای این قابلیت، ابتدا acceptance criteria بنویس، سپس migration/schema، service، API، UI، audit، test و documentation را در همان PR انجام بده. & Unit + integration + e2e؛ تست permission، validation، audit و حالت خطا. duplicate submission/vote و privacy/anonymity را تست کن. \\ \hline
175 & Professor evaluation & Survey/Question/Response/Answer را با anonymity و min response count کامل کن؛ dashboard results بساز؛ export امن اضافه کن؛ تست duplicate و privacy بنویس. برای این قابلیت، ابتدا acceptance criteria بنویس، سپس migration/schema، service، API، UI، audit، test و documentation را در همان PR انجام بده. & Unit + integration + e2e؛ تست permission، validation، audit و حالت خطا. duplicate submission/vote و privacy/anonymity را تست کن. \\ \hline
176 & course feedback & Survey/Question/Response/Answer را با anonymity و min response count کامل کن؛ dashboard results بساز؛ export امن اضافه کن؛ تست duplicate و privacy بنویس. برای این قابلیت، ابتدا acceptance criteria بنویس، سپس migration/schema، service، API، UI، audit، test و documentation را در همان PR انجام بده. & Unit + integration + e2e؛ تست permission، validation، audit و حالت خطا. duplicate submission/vote و privacy/anonymity را تست کن. \\ \hline
177 & session feedback & Survey/Question/Response/Answer را با anonymity و min response count کامل کن؛ dashboard results بساز؛ export امن اضافه کن؛ تست duplicate و privacy بنویس. برای این قابلیت، ابتدا acceptance criteria بنویس، سپس migration/schema، service، API، UI، audit، test و documentation را در همان PR انجام بده. & Unit + integration + e2e؛ تست permission، validation، audit و حالت خطا. duplicate submission/vote و privacy/anonymity را تست کن. \\ \hline
178 & تحلیل sentiment ساده & Survey/Question/Response/Answer را با anonymity و min response count کامل کن؛ dashboard results بساز؛ export امن اضافه کن؛ تست duplicate و privacy بنویس. برای این قابلیت، ابتدا acceptance criteria بنویس، سپس migration/schema، service، API، UI، audit، test و documentation را در همان PR انجام بده. & Unit + integration + e2e؛ تست permission، validation، audit و حالت خطا. duplicate submission/vote و privacy/anonymity را تست کن. \\ \hline
179 & keyword extraction & Survey/Question/Response/Answer را با anonymity و min response count کامل کن؛ dashboard results بساز؛ export امن اضافه کن؛ تست duplicate و privacy بنویس. برای این قابلیت، ابتدا acceptance criteria بنویس، سپس migration/schema، service، API، UI، audit، test و documentation را در همان PR انجام بده. & Unit + integration + e2e؛ تست permission، validation، audit و حالت خطا. duplicate submission/vote و privacy/anonymity را تست کن. \\ \hline
180 & comparison بین TAها & Survey/Question/Response/Answer را با anonymity و min response count کامل کن؛ dashboard results بساز؛ export امن اضافه کن؛ تست duplicate و privacy بنویس. برای این قابلیت، ابتدا acceptance criteria بنویس، سپس migration/schema، service، API، UI، audit، test و documentation را در همان PR انجام بده. & Unit + integration + e2e؛ تست permission، validation، audit و حالت خطا. duplicate submission/vote و privacy/anonymity را تست کن. \\ \hline
\end{longtable}
\subsection{K) Poll و زمان بندی}

\begin{longtable}{|>{\raggedleft\arraybackslash}p{0.065\textwidth}|>{\raggedleft\arraybackslash}p{0.21\textwidth}|>{\raggedleft\arraybackslash}p{0.42\textwidth}|>{\raggedleft\arraybackslash}p{0.18\textwidth}|}
\hline
\textbf{شماره} & \textbf{قابلیت} & \textbf{مراحل پیاده سازی} & \textbf{تست و پذیرش} \\\hline
\endfirsthead
\hline
\textbf{شماره} & \textbf{قابلیت} & \textbf{مراحل پیاده سازی} & \textbf{تست و پذیرش} \\\hline
\endhead
181 & poll انتخاب زمان کلاس & Poll/Option/Vote را با respondentHash و conflict detection کامل کن؛ UI vote/result بساز؛ reminder اضافه کن؛ تست duplicate vote و anonymity بنویس. برای این قابلیت، ابتدا acceptance criteria بنویس، سپس migration/schema، service، API، UI، audit، test و documentation را در همان PR انجام بده. & Unit + integration + e2e؛ تست permission، validation، audit و حالت خطا. duplicate submission/vote و privacy/anonymity را تست کن. \\ \hline
182 & poll انتخاب زمان حل تمرین & Poll/Option/Vote را با respondentHash و conflict detection کامل کن؛ UI vote/result بساز؛ reminder اضافه کن؛ تست duplicate vote و anonymity بنویس. برای این قابلیت، ابتدا acceptance criteria بنویس، سپس migration/schema، service، API، UI، audit، test و documentation را در همان PR انجام بده. & Unit + integration + e2e؛ تست permission، validation، audit و حالت خطا. duplicate submission/vote و privacy/anonymity را تست کن. \\ \hline
183 & poll جبرانی & Poll/Option/Vote را با respondentHash و conflict detection کامل کن؛ UI vote/result بساز؛ reminder اضافه کن؛ تست duplicate vote و anonymity بنویس. برای این قابلیت، ابتدا acceptance criteria بنویس، سپس migration/schema، service، API، UI، audit، test و documentation را در همان PR انجام بده. & Unit + integration + e2e؛ تست permission، validation، audit و حالت خطا. duplicate submission/vote و privacy/anonymity را تست کن. \\ \hline
184 & poll پروژه & Poll/Option/Vote را با respondentHash و conflict detection کامل کن؛ UI vote/result بساز؛ reminder اضافه کن؛ تست duplicate vote و anonymity بنویس. برای این قابلیت، ابتدا acceptance criteria بنویس، سپس migration/schema، service، API، UI، audit، test و documentation را در همان PR انجام بده. & Unit + integration + e2e؛ تست permission، validation، audit و حالت خطا. duplicate submission/vote و privacy/anonymity را تست کن. \\ \hline
185 & anonymous poll & Poll/Option/Vote را با respondentHash و conflict detection کامل کن؛ UI vote/result بساز؛ reminder اضافه کن؛ تست duplicate vote و anonymity بنویس. برای این قابلیت، ابتدا acceptance criteria بنویس، سپس migration/schema، service، API، UI، audit، test و documentation را در همان PR انجام بده. & Unit + integration + e2e؛ تست permission، validation، audit و حالت خطا. duplicate submission/vote و privacy/anonymity را تست کن. \\ \hline
186 & جلوگیری از vote duplicate & Poll/Option/Vote را با respondentHash و conflict detection کامل کن؛ UI vote/result بساز؛ reminder اضافه کن؛ تست duplicate vote و anonymity بنویس. برای این قابلیت، ابتدا acceptance criteria بنویس، سپس migration/schema، service، API، UI، audit، test و documentation را در همان PR انجام بده. & Unit + integration + e2e؛ تست permission، validation، audit و حالت خطا. duplicate submission/vote و privacy/anonymity را تست کن. \\ \hline
187 & multiple slot voting & Poll/Option/Vote را با respondentHash و conflict detection کامل کن؛ UI vote/result بساز؛ reminder اضافه کن؛ تست duplicate vote و anonymity بنویس. برای این قابلیت، ابتدا acceptance criteria بنویس، سپس migration/schema، service، API، UI، audit، test و documentation را در همان PR انجام بده. & Unit + integration + e2e؛ تست permission، validation، audit و حالت خطا. duplicate submission/vote و privacy/anonymity را تست کن. \\ \hline
188 & deadline poll & Poll/Option/Vote را با respondentHash و conflict detection کامل کن؛ UI vote/result بساز؛ reminder اضافه کن؛ تست duplicate vote و anonymity بنویس. برای این قابلیت، ابتدا acceptance criteria بنویس، سپس migration/schema، service، API، UI، audit، test و documentation را در همان PR انجام بده. & Unit + integration + e2e؛ تست permission، validation، audit و حالت خطا. duplicate submission/vote و privacy/anonymity را تست کن. \\ \hline
189 & close poll & Poll/Option/Vote را با respondentHash و conflict detection کامل کن؛ UI vote/result بساز؛ reminder اضافه کن؛ تست duplicate vote و anonymity بنویس. برای این قابلیت، ابتدا acceptance criteria بنویس، سپس migration/schema، service، API، UI، audit، test و documentation را در همان PR انجام بده. & Unit + integration + e2e؛ تست permission، validation، audit و حالت خطا. duplicate submission/vote و privacy/anonymity را تست کن. \\ \hline
190 & نمایش نتیجه بعد از بسته شدن & Poll/Option/Vote را با respondentHash و conflict detection کامل کن؛ UI vote/result بساز؛ reminder اضافه کن؛ تست duplicate vote و anonymity بنویس. برای این قابلیت، ابتدا acceptance criteria بنویس، سپس migration/schema، service، API، UI، audit، test و documentation را در همان PR انجام بده. & Unit + integration + e2e؛ تست permission، validation، audit و حالت خطا. duplicate submission/vote و privacy/anonymity را تست کن. \\ \hline
191 & export poll result & Poll/Option/Vote را با respondentHash و conflict detection کامل کن؛ UI vote/result بساز؛ reminder اضافه کن؛ تست duplicate vote و anonymity بنویس. برای این قابلیت، ابتدا acceptance criteria بنویس، سپس migration/schema، service، API، UI، audit، test و documentation را در همان PR انجام بده. & Unit + integration + e2e؛ تست permission، validation، audit و حالت خطا. duplicate submission/vote و privacy/anonymity را تست کن. \\ \hline
192 & notify participants & Poll/Option/Vote را با respondentHash و conflict detection کامل کن؛ UI vote/result بساز؛ reminder اضافه کن؛ تست duplicate vote و anonymity بنویس. برای این قابلیت، ابتدا acceptance criteria بنویس، سپس migration/schema، service، API، UI، audit، test و documentation را در همان PR انجام بده. & Unit + integration + e2e؛ تست permission، validation، audit و حالت خطا. duplicate submission/vote و privacy/anonymity را تست کن. \\ \hline
193 & calendar integration & Poll/Option/Vote را با respondentHash و conflict detection کامل کن؛ UI vote/result بساز؛ reminder اضافه کن؛ تست duplicate vote و anonymity بنویس. برای این قابلیت، ابتدا acceptance criteria بنویس، سپس migration/schema، service، API، UI، audit، test و documentation را در همان PR انجام بده. & Unit + integration + e2e؛ تست permission، validation، audit و حالت خطا. duplicate submission/vote و privacy/anonymity را تست کن. \\ \hline
194 & conflict detection & Poll/Option/Vote را با respondentHash و conflict detection کامل کن؛ UI vote/result بساز؛ reminder اضافه کن؛ تست duplicate vote و anonymity بنویس. برای این قابلیت، ابتدا acceptance criteria بنویس، سپس migration/schema، service، API، UI، audit، test و documentation را در همان PR انجام بده. & Unit + integration + e2e؛ تست permission، validation، audit و حالت خطا. duplicate submission/vote و privacy/anonymity را تست کن. \\ \hline
195 & recommendation بهترین زمان & Poll/Option/Vote را با respondentHash و conflict detection کامل کن؛ UI vote/result بساز؛ reminder اضافه کن؛ تست duplicate vote و anonymity بنویس. برای این قابلیت، ابتدا acceptance criteria بنویس، سپس migration/schema، service، API، UI، audit، test و documentation را در همان PR انجام بده. & Unit + integration + e2e؛ تست permission، validation، audit و حالت خطا. duplicate submission/vote و privacy/anonymity را تست کن. \\ \hline
196 & poll reminder & Poll/Option/Vote را با respondentHash و conflict detection کامل کن؛ UI vote/result بساز؛ reminder اضافه کن؛ تست duplicate vote و anonymity بنویس. برای این قابلیت، ابتدا acceptance criteria بنویس، سپس migration/schema، service، API، UI، audit، test و documentation را در همان PR انجام بده. & Unit + integration + e2e؛ تست permission، validation، audit و حالت خطا. duplicate submission/vote و privacy/anonymity را تست کن. \\ \hline
197 & تغییر رأی تا قبل deadline & Poll/Option/Vote را با respondentHash و conflict detection کامل کن؛ UI vote/result بساز؛ reminder اضافه کن؛ تست duplicate vote و anonymity بنویس. برای این قابلیت، ابتدا acceptance criteria بنویس، سپس migration/schema، service، API، UI، audit، test و documentation را در همان PR انجام بده. & Unit + integration + e2e؛ تست permission، validation، audit و حالت خطا. duplicate submission/vote و privacy/anonymity را تست کن. \\ \hline
198 & محدودیت رأی برای enrolled students & Poll/Option/Vote را با respondentHash و conflict detection کامل کن؛ UI vote/result بساز؛ reminder اضافه کن؛ تست duplicate vote و anonymity بنویس. برای این قابلیت، ابتدا acceptance criteria بنویس، سپس migration/schema، service، API، UI، audit، test و documentation را در همان PR انجام بده. & Unit + integration + e2e؛ تست permission، validation، audit و حالت خطا. duplicate submission/vote و privacy/anonymity را تست کن. \\ \hline
199 & audit vote & Poll/Option/Vote را با respondentHash و conflict detection کامل کن؛ UI vote/result بساز؛ reminder اضافه کن؛ تست duplicate vote و anonymity بنویس. برای این قابلیت، ابتدا acceptance criteria بنویس، سپس migration/schema، service، API، UI، audit، test و documentation را در همان PR انجام بده. & Unit + integration + e2e؛ تست permission، validation، audit و حالت خطا. duplicate submission/vote و privacy/anonymity را تست کن. \\ \hline
200 & privacy test برای poll & Poll/Option/Vote را با respondentHash و conflict detection کامل کن؛ UI vote/result بساز؛ reminder اضافه کن؛ تست duplicate vote و anonymity بنویس. برای این قابلیت، ابتدا acceptance criteria بنویس، سپس migration/schema، service، API، UI، audit، test و documentation را در همان PR انجام بده. & Unit + integration + e2e؛ تست permission، validation، audit و حالت خطا. duplicate submission/vote و privacy/anonymity را تست کن. \\ \hline
\end{longtable}
\subsection{L) گواهی TA/Head TA}

\begin{longtable}{|>{\raggedleft\arraybackslash}p{0.065\textwidth}|>{\raggedleft\arraybackslash}p{0.21\textwidth}|>{\raggedleft\arraybackslash}p{0.42\textwidth}|>{\raggedleft\arraybackslash}p{0.18\textwidth}|}
\hline
\textbf{شماره} & \textbf{قابلیت} & \textbf{مراحل پیاده سازی} & \textbf{تست و پذیرش} \\\hline
\endfirsthead
\hline
\textbf{شماره} & \textbf{قابلیت} & \textbf{مراحل پیاده سازی} & \textbf{تست و پذیرش} \\\hline
\endhead
201 & درخواست گواهی توسط TA & CertificateRequest/Certificate را با approval/issue/revoke کامل کن؛ PDF/QR/verify page بساز؛ bulk issue اضافه کن؛ تست eligibility، revoke و public privacy بنویس. برای این قابلیت، ابتدا acceptance criteria بنویس، سپس migration/schema، service، API، UI، audit، test و documentation را در همان PR انجام بده. & Unit + integration + e2e؛ تست permission، validation، audit و حالت خطا. signed URL، public verification و revoke/delete را تست کن. \\ \hline
202 & eligibility check & CertificateRequest/Certificate را با approval/issue/revoke کامل کن؛ PDF/QR/verify page بساز؛ bulk issue اضافه کن؛ تست eligibility، revoke و public privacy بنویس. برای این قابلیت، ابتدا acceptance criteria بنویس، سپس migration/schema، service، API، UI، audit، test و documentation را در همان PR انجام بده. & Unit + integration + e2e؛ تست permission، validation، audit و حالت خطا. signed URL، public verification و revoke/delete را تست کن. \\ \hline
203 & professor approval & CertificateRequest/Certificate را با approval/issue/revoke کامل کن؛ PDF/QR/verify page بساز؛ bulk issue اضافه کن؛ تست eligibility، revoke و public privacy بنویس. برای این قابلیت، ابتدا acceptance criteria بنویس، سپس migration/schema، service، API، UI، audit، test و documentation را در همان PR انجام بده. & Unit + integration + e2e؛ تست permission، validation، audit و حالت خطا. signed URL، public verification و revoke/delete را تست کن. \\ \hline
204 & education approval & CertificateRequest/Certificate را با approval/issue/revoke کامل کن؛ PDF/QR/verify page بساز؛ bulk issue اضافه کن؛ تست eligibility، revoke و public privacy بنویس. برای این قابلیت، ابتدا acceptance criteria بنویس، سپس migration/schema، service، API، UI، audit، test و documentation را در همان PR انجام بده. & Unit + integration + e2e؛ تست permission، validation، audit و حالت خطا. signed URL، public verification و revoke/delete را تست کن. \\ \hline
205 & issue certificate & CertificateRequest/Certificate را با approval/issue/revoke کامل کن؛ PDF/QR/verify page بساز؛ bulk issue اضافه کن؛ تست eligibility، revoke و public privacy بنویس. برای این قابلیت، ابتدا acceptance criteria بنویس، سپس migration/schema، service، API، UI، audit، test و documentation را در همان PR انجام بده. & Unit + integration + e2e؛ تست permission، validation، audit و حالت خطا. signed URL، public verification و revoke/delete را تست کن. \\ \hline
206 & PDF certificate generation & CertificateRequest/Certificate را با approval/issue/revoke کامل کن؛ PDF/QR/verify page بساز؛ bulk issue اضافه کن؛ تست eligibility، revoke و public privacy بنویس. برای این قابلیت، ابتدا acceptance criteria بنویس، سپس migration/schema، service، API، UI، audit، test و documentation را در همان PR انجام بده. & Unit + integration + e2e؛ تست permission، validation، audit و حالت خطا. signed URL، public verification و revoke/delete را تست کن. \\ \hline
207 & QR code & CertificateRequest/Certificate را با approval/issue/revoke کامل کن؛ PDF/QR/verify page بساز؛ bulk issue اضافه کن؛ تست eligibility، revoke و public privacy بنویس. برای این قابلیت، ابتدا acceptance criteria بنویس، سپس migration/schema، service، API، UI، audit، test و documentation را در همان PR انجام بده. & Unit + integration + e2e؛ تست permission، validation، audit و حالت خطا. signed URL، public verification و revoke/delete را تست کن. \\ \hline
208 & public verification page & CertificateRequest/Certificate را با approval/issue/revoke کامل کن؛ PDF/QR/verify page بساز؛ bulk issue اضافه کن؛ تست eligibility، revoke و public privacy بنویس. برای این قابلیت، ابتدا acceptance criteria بنویس، سپس migration/schema، service، API، UI، audit، test و documentation را در همان PR انجام بده. & Unit + integration + e2e؛ تست permission، validation، audit و حالت خطا. signed URL، public verification و revoke/delete را تست کن. \\ \hline
209 & tracking code & CertificateRequest/Certificate را با approval/issue/revoke کامل کن؛ PDF/QR/verify page بساز؛ bulk issue اضافه کن؛ تست eligibility، revoke و public privacy بنویس. برای این قابلیت، ابتدا acceptance criteria بنویس، سپس migration/schema، service، API، UI، audit، test و documentation را در همان PR انجام بده. & Unit + integration + e2e؛ تست permission، validation، audit و حالت خطا. signed URL، public verification و revoke/delete را تست کن. \\ \hline
210 & verification token hash & CertificateRequest/Certificate را با approval/issue/revoke کامل کن؛ PDF/QR/verify page بساز؛ bulk issue اضافه کن؛ تست eligibility، revoke و public privacy بنویس. برای این قابلیت، ابتدا acceptance criteria بنویس، سپس migration/schema، service، API، UI، audit، test و documentation را در همان PR انجام بده. & Unit + integration + e2e؛ تست permission، validation، audit و حالت خطا. signed URL، public verification و revoke/delete را تست کن. \\ \hline
211 & revoke certificate & CertificateRequest/Certificate را با approval/issue/revoke کامل کن؛ PDF/QR/verify page بساز؛ bulk issue اضافه کن؛ تست eligibility، revoke و public privacy بنویس. برای این قابلیت، ابتدا acceptance criteria بنویس، سپس migration/schema، service، API، UI، audit، test و documentation را در همان PR انجام بده. & Unit + integration + e2e؛ تست permission، validation، audit و حالت خطا. signed URL، public verification و revoke/delete را تست کن. \\ \hline
212 & revoke reason & CertificateRequest/Certificate را با approval/issue/revoke کامل کن؛ PDF/QR/verify page بساز؛ bulk issue اضافه کن؛ تست eligibility، revoke و public privacy بنویس. برای این قابلیت، ابتدا acceptance criteria بنویس، سپس migration/schema، service، API، UI، audit، test و documentation را در همان PR انجام بده. & Unit + integration + e2e؛ تست permission، validation، audit و حالت خطا. signed URL، public verification و revoke/delete را تست کن. \\ \hline
213 & bulk issue certificates & CertificateRequest/Certificate را با approval/issue/revoke کامل کن؛ PDF/QR/verify page بساز؛ bulk issue اضافه کن؛ تست eligibility، revoke و public privacy بنویس. برای این قابلیت، ابتدا acceptance criteria بنویس، سپس migration/schema، service، API، UI، audit، test و documentation را در همان PR انجام بده. & Unit + integration + e2e؛ تست permission، validation، audit و حالت خطا. signed URL، public verification و revoke/delete را تست کن. \\ \hline
214 & certificate template management & CertificateRequest/Certificate را با approval/issue/revoke کامل کن؛ PDF/QR/verify page بساز؛ bulk issue اضافه کن؛ تست eligibility، revoke و public privacy بنویس. برای این قابلیت، ابتدا acceptance criteria بنویس، سپس migration/schema، service، API، UI، audit، test و documentation را در همان PR انجام بده. & Unit + integration + e2e؛ تست permission، validation، audit و حالت خطا. signed URL، public verification و revoke/delete را تست کن. \\ \hline
215 & فارسی سازی PDF & CertificateRequest/Certificate را با approval/issue/revoke کامل کن؛ PDF/QR/verify page بساز؛ bulk issue اضافه کن؛ تست eligibility، revoke و public privacy بنویس. برای این قابلیت، ابتدا acceptance criteria بنویس، سپس migration/schema، service، API، UI، audit، test و documentation را در همان PR انجام بده. & Unit + integration + e2e؛ تست permission، validation، audit و حالت خطا. signed URL، public verification و revoke/delete را تست کن. \\ \hline
216 & خروجی قابل چاپ & CertificateRequest/Certificate را با approval/issue/revoke کامل کن؛ PDF/QR/verify page بساز؛ bulk issue اضافه کن؛ تست eligibility، revoke و public privacy بنویس. برای این قابلیت، ابتدا acceptance criteria بنویس، سپس migration/schema، service، API، UI، audit، test و documentation را در همان PR انجام بده. & Unit + integration + e2e؛ تست permission، validation، audit و حالت خطا. signed URL، public verification و revoke/delete را تست کن. \\ \hline
217 & امضای دیجیتال اختیاری & CertificateRequest/Certificate را با approval/issue/revoke کامل کن؛ PDF/QR/verify page بساز؛ bulk issue اضافه کن؛ تست eligibility، revoke و public privacy بنویس. برای این قابلیت، ابتدا acceptance criteria بنویس، سپس migration/schema، service، API، UI، audit، test و documentation را در همان PR انجام بده. & Unit + integration + e2e؛ تست permission، validation، audit و حالت خطا. signed URL، public verification و revoke/delete را تست کن. \\ \hline
218 & شماره گواهی رسمی & CertificateRequest/Certificate را با approval/issue/revoke کامل کن؛ PDF/QR/verify page بساز؛ bulk issue اضافه کن؛ تست eligibility، revoke و public privacy بنویس. برای این قابلیت، ابتدا acceptance criteria بنویس، سپس migration/schema، service، API، UI، audit، test و documentation را در همان PR انجام بده. & Unit + integration + e2e؛ تست permission، validation، audit و حالت خطا. signed URL، public verification و revoke/delete را تست کن. \\ \hline
219 & certificate audit & CertificateRequest/Certificate را با approval/issue/revoke کامل کن؛ PDF/QR/verify page بساز؛ bulk issue اضافه کن؛ تست eligibility، revoke و public privacy بنویس. برای این قابلیت، ابتدا acceptance criteria بنویس، سپس migration/schema، service، API، UI، audit، test و documentation را در همان PR انجام بده. & Unit + integration + e2e؛ تست permission، validation، audit و حالت خطا. signed URL، public verification و revoke/delete را تست کن. \\ \hline
220 & certificate export & CertificateRequest/Certificate را با approval/issue/revoke کامل کن؛ PDF/QR/verify page بساز؛ bulk issue اضافه کن؛ تست eligibility، revoke و public privacy بنویس. برای این قابلیت، ابتدا acceptance criteria بنویس، سپس migration/schema، service، API، UI، audit، test و documentation را در همان PR انجام بده. & Unit + integration + e2e؛ تست permission، validation، audit و حالت خطا. signed URL، public verification و revoke/delete را تست کن. \\ \hline
\end{longtable}
\subsection{M) فایل ها و Storage}

\begin{longtable}{|>{\raggedleft\arraybackslash}p{0.065\textwidth}|>{\raggedleft\arraybackslash}p{0.21\textwidth}|>{\raggedleft\arraybackslash}p{0.42\textwidth}|>{\raggedleft\arraybackslash}p{0.18\textwidth}|}
\hline
\textbf{شماره} & \textbf{قابلیت} & \textbf{مراحل پیاده سازی} & \textbf{تست و پذیرش} \\\hline
\endfirsthead
\hline
\textbf{شماره} & \textbf{قابلیت} & \textbf{مراحل پیاده سازی} & \textbf{تست و پذیرش} \\\hline
\endhead
221 & upload رزومه & UploadedFile و storage adapter را کامل کن؛ signed URL، checksum، scanStatus و audit اضافه کن؛ UI preview بساز؛ تست permission و orphan cleanup بنویس. برای این قابلیت، ابتدا acceptance criteria بنویس، سپس migration/schema، service، API، UI، audit، test و documentation را در همان PR انجام بده. & Unit + integration + e2e؛ تست permission، validation، audit و حالت خطا. signed URL، public verification و revoke/delete را تست کن. \\ \hline
222 & upload فایل درس & UploadedFile و storage adapter را کامل کن؛ signed URL، checksum، scanStatus و audit اضافه کن؛ UI preview بساز؛ تست permission و orphan cleanup بنویس. برای این قابلیت، ابتدا acceptance criteria بنویس، سپس migration/schema، service، API، UI، audit، test و documentation را در همان PR انجام بده. & Unit + integration + e2e؛ تست permission، validation، audit و حالت خطا. signed URL، public verification و revoke/delete را تست کن. \\ \hline
223 & upload فایل تکلیف & UploadedFile و storage adapter را کامل کن؛ signed URL، checksum، scanStatus و audit اضافه کن؛ UI preview بساز؛ تست permission و orphan cleanup بنویس. برای این قابلیت، ابتدا acceptance criteria بنویس، سپس migration/schema، service، API، UI، audit، test و documentation را در همان PR انجام بده. & Unit + integration + e2e؛ تست permission، validation، audit و حالت خطا. signed URL، public verification و revoke/delete را تست کن. \\ \hline
224 & S3/MinIO storage & UploadedFile و storage adapter را کامل کن؛ signed URL، checksum، scanStatus و audit اضافه کن؛ UI preview بساز؛ تست permission و orphan cleanup بنویس. برای این قابلیت، ابتدا acceptance criteria بنویس، سپس migration/schema، service، API، UI، audit، test و documentation را در همان PR انجام بده. & Unit + integration + e2e؛ تست permission، validation، audit و حالت خطا. signed URL، public verification و revoke/delete را تست کن. \\ \hline
225 & signed URL & UploadedFile و storage adapter را کامل کن؛ signed URL، checksum، scanStatus و audit اضافه کن؛ UI preview بساز؛ تست permission و orphan cleanup بنویس. برای این قابلیت، ابتدا acceptance criteria بنویس، سپس migration/schema، service، API، UI، audit، test و documentation را در همان PR انجام بده. & Unit + integration + e2e؛ تست permission، validation، audit و حالت خطا. signed URL، public verification و revoke/delete را تست کن. \\ \hline
226 & private file access & UploadedFile و storage adapter را کامل کن؛ signed URL، checksum، scanStatus و audit اضافه کن؛ UI preview بساز؛ تست permission و orphan cleanup بنویس. برای این قابلیت، ابتدا acceptance criteria بنویس، سپس migration/schema، service، API، UI، audit، test و documentation را در همان PR انجام بده. & Unit + integration + e2e؛ تست permission، validation، audit و حالت خطا. signed URL، public verification و revoke/delete را تست کن. \\ \hline
227 & visibility level & UploadedFile و storage adapter را کامل کن؛ signed URL، checksum، scanStatus و audit اضافه کن؛ UI preview بساز؛ تست permission و orphan cleanup بنویس. برای این قابلیت، ابتدا acceptance criteria بنویس، سپس migration/schema، service، API، UI، audit، test و documentation را در همان PR انجام بده. & Unit + integration + e2e؛ تست permission، validation، audit و حالت خطا. signed URL، public verification و revoke/delete را تست کن. \\ \hline
228 & virus scan status & UploadedFile و storage adapter را کامل کن؛ signed URL، checksum، scanStatus و audit اضافه کن؛ UI preview بساز؛ تست permission و orphan cleanup بنویس. برای این قابلیت، ابتدا acceptance criteria بنویس، سپس migration/schema، service، API، UI، audit، test و documentation را در همان PR انجام بده. & Unit + integration + e2e؛ تست permission، validation، audit و حالت خطا. signed URL، public verification و revoke/delete را تست کن. \\ \hline
229 & checksum SHA256 & UploadedFile و storage adapter را کامل کن؛ signed URL، checksum، scanStatus و audit اضافه کن؛ UI preview بساز؛ تست permission و orphan cleanup بنویس. برای این قابلیت، ابتدا acceptance criteria بنویس، سپس migration/schema، service، API، UI، audit، test و documentation را در همان PR انجام بده. & Unit + integration + e2e؛ تست permission، validation، audit و حالت خطا. signed URL، public verification و revoke/delete را تست کن. \\ \hline
230 & file ownership & UploadedFile و storage adapter را کامل کن؛ signed URL، checksum، scanStatus و audit اضافه کن؛ UI preview بساز؛ تست permission و orphan cleanup بنویس. برای این قابلیت، ابتدا acceptance criteria بنویس، سپس migration/schema، service، API، UI، audit، test و documentation را در همان PR انجام بده. & Unit + integration + e2e؛ تست permission، validation، audit و حالت خطا. signed URL، public verification و revoke/delete را تست کن. \\ \hline
231 & soft delete file & UploadedFile و storage adapter را کامل کن؛ signed URL، checksum، scanStatus و audit اضافه کن؛ UI preview بساز؛ تست permission و orphan cleanup بنویس. برای این قابلیت، ابتدا acceptance criteria بنویس، سپس migration/schema، service، API، UI، audit، test و documentation را در همان PR انجام بده. & Unit + integration + e2e؛ تست permission، validation، audit و حالت خطا. signed URL، public verification و revoke/delete را تست کن. \\ \hline
232 & audit download & UploadedFile و storage adapter را کامل کن؛ signed URL، checksum، scanStatus و audit اضافه کن؛ UI preview بساز؛ تست permission و orphan cleanup بنویس. برای این قابلیت، ابتدا acceptance criteria بنویس، سپس migration/schema، service، API، UI، audit، test و documentation را در همان PR انجام بده. & Unit + integration + e2e؛ تست permission، validation، audit و حالت خطا. signed URL، public verification و revoke/delete را تست کن. \\ \hline
233 & file preview & UploadedFile و storage adapter را کامل کن؛ signed URL، checksum، scanStatus و audit اضافه کن؛ UI preview بساز؛ تست permission و orphan cleanup بنویس. برای این قابلیت، ابتدا acceptance criteria بنویس، سپس migration/schema، service، API، UI، audit، test و documentation را در همان PR انجام بده. & Unit + integration + e2e؛ تست permission، validation، audit و حالت خطا. signed URL، public verification و revoke/delete را تست کن. \\ \hline
234 & PDF preview & UploadedFile و storage adapter را کامل کن؛ signed URL، checksum، scanStatus و audit اضافه کن؛ UI preview بساز؛ تست permission و orphan cleanup بنویس. برای این قابلیت، ابتدا acceptance criteria بنویس، سپس migration/schema، service، API، UI، audit، test و documentation را در همان PR انجام بده. & Unit + integration + e2e؛ تست permission، validation، audit و حالت خطا. signed URL، public verification و revoke/delete را تست کن. \\ \hline
235 & image preview & UploadedFile و storage adapter را کامل کن؛ signed URL، checksum، scanStatus و audit اضافه کن؛ UI preview بساز؛ تست permission و orphan cleanup بنویس. برای این قابلیت، ابتدا acceptance criteria بنویس، سپس migration/schema، service، API، UI، audit، test و documentation را در همان PR انجام بده. & Unit + integration + e2e؛ تست permission، validation، audit و حالت خطا. signed URL، public verification و revoke/delete را تست کن. \\ \hline
236 & attachment to message & UploadedFile و storage adapter را کامل کن؛ signed URL، checksum، scanStatus و audit اضافه کن؛ UI preview بساز؛ تست permission و orphan cleanup بنویس. برای این قابلیت، ابتدا acceptance criteria بنویس، سپس migration/schema، service، API، UI، audit، test و documentation را در همان PR انجام بده. & Unit + integration + e2e؛ تست permission، validation، audit و حالت خطا. signed URL، public verification و revoke/delete را تست کن. \\ \hline
237 & attachment to submission & UploadedFile و storage adapter را کامل کن؛ signed URL، checksum، scanStatus و audit اضافه کن؛ UI preview بساز؛ تست permission و orphan cleanup بنویس. برای این قابلیت، ابتدا acceptance criteria بنویس، سپس migration/schema، service، API، UI، audit، test و documentation را در همان PR انجام بده. & Unit + integration + e2e؛ تست permission، validation، audit و حالت خطا. signed URL، public verification و revoke/delete را تست کن. \\ \hline
238 & attachment to certificate & UploadedFile و storage adapter را کامل کن؛ signed URL، checksum، scanStatus و audit اضافه کن؛ UI preview بساز؛ تست permission و orphan cleanup بنویس. برای این قابلیت، ابتدا acceptance criteria بنویس، سپس migration/schema، service، API، UI، audit، test و documentation را در همان PR انجام بده. & Unit + integration + e2e؛ تست permission، validation، audit و حالت خطا. signed URL، public verification و revoke/delete را تست کن. \\ \hline
239 & storage quota & UploadedFile و storage adapter را کامل کن؛ signed URL، checksum، scanStatus و audit اضافه کن؛ UI preview بساز؛ تست permission و orphan cleanup بنویس. برای این قابلیت، ابتدا acceptance criteria بنویس، سپس migration/schema، service، API، UI، audit، test و documentation را در همان PR انجام بده. & Unit + integration + e2e؛ تست permission، validation، audit و حالت خطا. signed URL، public verification و revoke/delete را تست کن. \\ \hline
240 & cleanup orphan files & UploadedFile و storage adapter را کامل کن؛ signed URL، checksum، scanStatus و audit اضافه کن؛ UI preview بساز؛ تست permission و orphan cleanup بنویس. برای این قابلیت، ابتدا acceptance criteria بنویس، سپس migration/schema، service، API، UI، audit، test و documentation را در همان PR انجام بده. & Unit + integration + e2e؛ تست permission، validation، audit و حالت خطا. signed URL، public verification و revoke/delete را تست کن. \\ \hline
\end{longtable}
\subsection{N) گزارش ها و Analytics}

\begin{longtable}{|>{\raggedleft\arraybackslash}p{0.065\textwidth}|>{\raggedleft\arraybackslash}p{0.21\textwidth}|>{\raggedleft\arraybackslash}p{0.42\textwidth}|>{\raggedleft\arraybackslash}p{0.18\textwidth}|}
\hline
\textbf{شماره} & \textbf{قابلیت} & \textbf{مراحل پیاده سازی} & \textbf{تست و پذیرش} \\\hline
\endfirsthead
\hline
\textbf{شماره} & \textbf{قابلیت} & \textbf{مراحل پیاده سازی} & \textbf{تست و پذیرش} \\\hline
\endhead
241 & گزارش تعداد TAها & query/report service و aggregationها را بساز؛ chart/export/dashboard اضافه کن؛ cache مناسب بگذار؛ تست permission، accuracy و export بنویس. برای این قابلیت، ابتدا acceptance criteria بنویس، سپس migration/schema، service، API، UI، audit، test و documentation را در همان PR انجام بده. & Unit + integration + e2e؛ تست permission، validation، audit و حالت خطا. \\ \hline
242 & گزارش applicationها & query/report service و aggregationها را بساز؛ chart/export/dashboard اضافه کن؛ cache مناسب بگذار؛ تست permission، accuracy و export بنویس. برای این قابلیت، ابتدا acceptance criteria بنویس، سپس migration/schema، service، API، UI، audit، test و documentation را در همان PR انجام بده. & Unit + integration + e2e؛ تست permission، validation، audit و حالت خطا. \\ \hline
243 & گزارش قبولی/ردی & query/report service و aggregationها را بساز؛ chart/export/dashboard اضافه کن؛ cache مناسب بگذار؛ تست permission، accuracy و export بنویس. برای این قابلیت، ابتدا acceptance criteria بنویس، سپس migration/schema، service، API، UI، audit، test و documentation را در همان PR انجام بده. & Unit + integration + e2e؛ تست permission، validation، audit و حالت خطا. \\ \hline
244 & گزارش workload TA & query/report service و aggregationها را بساز؛ chart/export/dashboard اضافه کن؛ cache مناسب بگذار؛ تست permission، accuracy و export بنویس. برای این قابلیت، ابتدا acceptance criteria بنویس، سپس migration/schema، service، API، UI، audit، test و documentation را در همان PR انجام بده. & Unit + integration + e2e؛ تست permission، validation، audit و حالت خطا. \\ \hline
245 & گزارش جلسات برگزارشده & query/report service و aggregationها را بساز؛ chart/export/dashboard اضافه کن؛ cache مناسب بگذار؛ تست permission، accuracy و export بنویس. برای این قابلیت، ابتدا acceptance criteria بنویس، سپس migration/schema، service، API، UI، audit، test و documentation را در همان PR انجام بده. & Unit + integration + e2e؛ تست permission، validation، audit و حالت خطا. \\ \hline
246 & گزارش attendance & query/report service و aggregationها را بساز؛ chart/export/dashboard اضافه کن؛ cache مناسب بگذار؛ تست permission، accuracy و export بنویس. برای این قابلیت، ابتدا acceptance criteria بنویس، سپس migration/schema، service، API، UI، audit، test و documentation را در همان PR انجام بده. & Unit + integration + e2e؛ تست permission، validation، audit و حالت خطا. \\ \hline
247 & گزارش میانگین نمرات & query/report service و aggregationها را بساز؛ chart/export/dashboard اضافه کن؛ cache مناسب بگذار؛ تست permission، accuracy و export بنویس. برای این قابلیت، ابتدا acceptance criteria بنویس، سپس migration/schema، service، API، UI، audit، test و documentation را در همان PR انجام بده. & Unit + integration + e2e؛ تست permission، validation، audit و حالت خطا. \\ \hline
248 & گزارش سرعت تصحیح & query/report service و aggregationها را بساز؛ chart/export/dashboard اضافه کن؛ cache مناسب بگذار؛ تست permission، accuracy و export بنویس. برای این قابلیت، ابتدا acceptance criteria بنویس، سپس migration/schema، service، API، UI، audit، test و documentation را در همان PR انجام بده. & Unit + integration + e2e؛ تست permission، validation، audit و حالت خطا. \\ \hline
249 & گزارش پیام ها & query/report service و aggregationها را بساز؛ chart/export/dashboard اضافه کن؛ cache مناسب بگذار؛ تست permission، accuracy و export بنویس. برای این قابلیت، ابتدا acceptance criteria بنویس، سپس migration/schema، service، API، UI، audit، test و documentation را در همان PR انجام بده. & Unit + integration + e2e؛ تست permission، validation، audit و حالت خطا. \\ \hline
250 & گزارش survey satisfaction & query/report service و aggregationها را بساز؛ chart/export/dashboard اضافه کن؛ cache مناسب بگذار؛ تست permission، accuracy و export بنویس. برای این قابلیت، ابتدا acceptance criteria بنویس، سپس migration/schema، service، API، UI، audit، test و documentation را در همان PR انجام بده. & Unit + integration + e2e؛ تست permission، validation، audit و حالت خطا. \\ \hline
251 & گزارش professor dashboard & query/report service و aggregationها را بساز؛ chart/export/dashboard اضافه کن؛ cache مناسب بگذار؛ تست permission، accuracy و export بنویس. برای این قابلیت، ابتدا acceptance criteria بنویس، سپس migration/schema، service، API، UI، audit، test و documentation را در همان PR انجام بده. & Unit + integration + e2e؛ تست permission، validation، audit و حالت خطا. \\ \hline
252 & گزارش Head TA dashboard & query/report service و aggregationها را بساز؛ chart/export/dashboard اضافه کن؛ cache مناسب بگذار؛ تست permission، accuracy و export بنویس. برای این قابلیت، ابتدا acceptance criteria بنویس، سپس migration/schema، service، API، UI، audit، test و documentation را در همان PR انجام بده. & Unit + integration + e2e؛ تست permission، validation، audit و حالت خطا. \\ \hline
253 & گزارش admin dashboard & query/report service و aggregationها را بساز؛ chart/export/dashboard اضافه کن؛ cache مناسب بگذار؛ تست permission، accuracy و export بنویس. برای این قابلیت، ابتدا acceptance criteria بنویس، سپس migration/schema، service، API، UI، audit، test و documentation را در همان PR انجام بده. & Unit + integration + e2e؛ تست permission، validation، audit و حالت خطا. \\ \hline
254 & export reports & query/report service و aggregationها را بساز؛ chart/export/dashboard اضافه کن؛ cache مناسب بگذار؛ تست permission، accuracy و export بنویس. برای این قابلیت، ابتدا acceptance criteria بنویس، سپس migration/schema، service، API، UI، audit، test و documentation را در همان PR انجام بده. & Unit + integration + e2e؛ تست permission، validation، audit و حالت خطا. \\ \hline
255 & chart و نمودار & query/report service و aggregationها را بساز؛ chart/export/dashboard اضافه کن؛ cache مناسب بگذار؛ تست permission، accuracy و export بنویس. برای این قابلیت، ابتدا acceptance criteria بنویس، سپس migration/schema، service، API، UI، audit، test و documentation را در همان PR انجام بده. & Unit + integration + e2e؛ تست permission، validation، audit و حالت خطا. \\ \hline
256 & activity timeline & query/report service و aggregationها را بساز؛ chart/export/dashboard اضافه کن؛ cache مناسب بگذار؛ تست permission، accuracy و export بنویس. برای این قابلیت، ابتدا acceptance criteria بنویس، سپس migration/schema، service، API، UI، audit، test و documentation را در همان PR انجام بده. & Unit + integration + e2e؛ تست permission، validation، audit و حالت خطا. \\ \hline
257 & audit timeline & query/report service و aggregationها را بساز؛ chart/export/dashboard اضافه کن؛ cache مناسب بگذار؛ تست permission، accuracy و export بنویس. برای این قابلیت، ابتدا acceptance criteria بنویس، سپس migration/schema، service، API، UI، audit، test و documentation را در همان PR انجام بده. & Unit + integration + e2e؛ تست permission، validation، audit و حالت خطا. \\ \hline
258 & risk flags & query/report service و aggregationها را بساز؛ chart/export/dashboard اضافه کن؛ cache مناسب بگذار؛ تست permission، accuracy و export بنویس. برای این قابلیت، ابتدا acceptance criteria بنویس، سپس migration/schema، service، API، UI، audit، test و documentation را در همان PR انجام بده. & Unit + integration + e2e؛ تست permission، validation، audit و حالت خطا. \\ \hline
259 & inactive TA detection & query/report service و aggregationها را بساز؛ chart/export/dashboard اضافه کن؛ cache مناسب بگذار؛ تست permission، accuracy و export بنویس. برای این قابلیت، ابتدا acceptance criteria بنویس، سپس migration/schema، service، API، UI، audit، test و documentation را در همان PR انجام بده. & Unit + integration + e2e؛ تست permission، validation، audit و حالت خطا. \\ \hline
260 & overworked TA detection & query/report service و aggregationها را بساز؛ chart/export/dashboard اضافه کن؛ cache مناسب بگذار؛ تست permission، accuracy و export بنویس. برای این قابلیت، ابتدا acceptance criteria بنویس، سپس migration/schema، service، API، UI، audit، test و documentation را در همان PR انجام بده. & Unit + integration + e2e؛ تست permission، validation، audit و حالت خطا. \\ \hline
\end{longtable}
\subsection{O) پنل ها و UI}

\begin{longtable}{|>{\raggedleft\arraybackslash}p{0.065\textwidth}|>{\raggedleft\arraybackslash}p{0.21\textwidth}|>{\raggedleft\arraybackslash}p{0.42\textwidth}|>{\raggedleft\arraybackslash}p{0.18\textwidth}|}
\hline
\textbf{شماره} & \textbf{قابلیت} & \textbf{مراحل پیاده سازی} & \textbf{تست و پذیرش} \\\hline
\endfirsthead
\hline
\textbf{شماره} & \textbf{قابلیت} & \textbf{مراحل پیاده سازی} & \textbf{تست و پذیرش} \\\hline
\endhead
261 & landing page مدرن & design system را با tokens، RTL، dark/light، responsive و animation کنترل شده بازسازی کن؛ صفحات role-based را polish کن؛ visual/e2e/a11y تست بنویس. برای این قابلیت، ابتدا acceptance criteria بنویس، سپس migration/schema، service، API، UI، audit، test و documentation را در همان PR انجام بده. & Unit + integration + e2e؛ تست permission، validation، audit و حالت خطا. \\ \hline
262 & hero animated & design system را با tokens، RTL، dark/light، responsive و animation کنترل شده بازسازی کن؛ صفحات role-based را polish کن؛ visual/e2e/a11y تست بنویس. برای این قابلیت، ابتدا acceptance criteria بنویس، سپس migration/schema، service، API، UI، audit، test و documentation را در همان PR انجام بده. & Unit + integration + e2e؛ تست permission، validation، audit و حالت خطا. \\ \hline
263 & bento grid features & design system را با tokens، RTL، dark/light، responsive و animation کنترل شده بازسازی کن؛ صفحات role-based را polish کن؛ visual/e2e/a11y تست بنویس. برای این قابلیت، ابتدا acceptance criteria بنویس، سپس migration/schema، service، API، UI، audit، test و documentation را در همان PR انجام بده. & Unit + integration + e2e؛ تست permission، validation، audit و حالت خطا. \\ \hline
264 & dashboard mockup & design system را با tokens، RTL، dark/light، responsive و animation کنترل شده بازسازی کن؛ صفحات role-based را polish کن؛ visual/e2e/a11y تست بنویس. برای این قابلیت، ابتدا acceptance criteria بنویس، سپس migration/schema، service، API، UI، audit، test و documentation را در همان PR انجام بده. & Unit + integration + e2e؛ تست permission، validation، audit و حالت خطا. \\ \hline
265 & role-based dashboard & design system را با tokens، RTL، dark/light، responsive و animation کنترل شده بازسازی کن؛ صفحات role-based را polish کن؛ visual/e2e/a11y تست بنویس. برای این قابلیت، ابتدا acceptance criteria بنویس، سپس migration/schema، service، API، UI، audit، test و documentation را در همان PR انجام بده. & Unit + integration + e2e؛ تست permission، validation، audit و حالت خطا. \\ \hline
266 & student dashboard & design system را با tokens، RTL، dark/light، responsive و animation کنترل شده بازسازی کن؛ صفحات role-based را polish کن؛ visual/e2e/a11y تست بنویس. برای این قابلیت، ابتدا acceptance criteria بنویس، سپس migration/schema، service، API، UI، audit، test و documentation را در همان PR انجام بده. & Unit + integration + e2e؛ تست permission، validation، audit و حالت خطا. \\ \hline
267 & TA dashboard & design system را با tokens، RTL، dark/light، responsive و animation کنترل شده بازسازی کن؛ صفحات role-based را polish کن؛ visual/e2e/a11y تست بنویس. برای این قابلیت، ابتدا acceptance criteria بنویس، سپس migration/schema، service، API، UI، audit، test و documentation را در همان PR انجام بده. & Unit + integration + e2e؛ تست permission، validation، audit و حالت خطا. \\ \hline
268 & Head TA dashboard & design system را با tokens، RTL، dark/light، responsive و animation کنترل شده بازسازی کن؛ صفحات role-based را polish کن؛ visual/e2e/a11y تست بنویس. برای این قابلیت، ابتدا acceptance criteria بنویس، سپس migration/schema، service، API، UI، audit، test و documentation را در همان PR انجام بده. & Unit + integration + e2e؛ تست permission، validation، audit و حالت خطا. \\ \hline
269 & Professor dashboard & design system را با tokens، RTL، dark/light، responsive و animation کنترل شده بازسازی کن؛ صفحات role-based را polish کن؛ visual/e2e/a11y تست بنویس. برای این قابلیت، ابتدا acceptance criteria بنویس، سپس migration/schema، service، API، UI، audit، test و documentation را در همان PR انجام بده. & Unit + integration + e2e؛ تست permission، validation، audit و حالت خطا. \\ \hline
270 & Admin dashboard & design system را با tokens، RTL، dark/light، responsive و animation کنترل شده بازسازی کن؛ صفحات role-based را polish کن؛ visual/e2e/a11y تست بنویس. برای این قابلیت، ابتدا acceptance criteria بنویس، سپس migration/schema، service، API، UI، audit، test و documentation را در همان PR انجام بده. & Unit + integration + e2e؛ تست permission، validation، audit و حالت خطا. \\ \hline
271 & dark mode & design system را با tokens، RTL، dark/light، responsive و animation کنترل شده بازسازی کن؛ صفحات role-based را polish کن؛ visual/e2e/a11y تست بنویس. برای این قابلیت، ابتدا acceptance criteria بنویس، سپس migration/schema، service، API، UI، audit، test و documentation را در همان PR انجام بده. & Unit + integration + e2e؛ تست permission، validation، audit و حالت خطا. \\ \hline
272 & light mode & design system را با tokens، RTL، dark/light، responsive و animation کنترل شده بازسازی کن؛ صفحات role-based را polish کن؛ visual/e2e/a11y تست بنویس. برای این قابلیت، ابتدا acceptance criteria بنویس، سپس migration/schema، service، API، UI، audit، test و documentation را در همان PR انجام بده. & Unit + integration + e2e؛ تست permission، validation، audit و حالت خطا. \\ \hline
273 & RTL کامل & design system را با tokens، RTL، dark/light، responsive و animation کنترل شده بازسازی کن؛ صفحات role-based را polish کن؛ visual/e2e/a11y تست بنویس. برای این قابلیت، ابتدا acceptance criteria بنویس، سپس migration/schema، service، API، UI، audit، test و documentation را در همان PR انجام بده. & Unit + integration + e2e؛ تست permission، validation، audit و حالت خطا. \\ \hline
274 & mobile responsive & design system را با tokens، RTL، dark/light، responsive و animation کنترل شده بازسازی کن؛ صفحات role-based را polish کن؛ visual/e2e/a11y تست بنویس. برای این قابلیت، ابتدا acceptance criteria بنویس، سپس migration/schema، service، API، UI، audit، test و documentation را در همان PR انجام بده. & Unit + integration + e2e؛ تست permission، validation، audit و حالت خطا. \\ \hline
275 & table responsive & design system را با tokens، RTL، dark/light، responsive و animation کنترل شده بازسازی کن؛ صفحات role-based را polish کن؛ visual/e2e/a11y تست بنویس. برای این قابلیت، ابتدا acceptance criteria بنویس، سپس migration/schema، service، API، UI، audit، test و documentation را در همان PR انجام بده. & Unit + integration + e2e؛ تست permission، validation، audit و حالت خطا. \\ \hline
276 & skeleton loading & design system را با tokens، RTL، dark/light، responsive و animation کنترل شده بازسازی کن؛ صفحات role-based را polish کن؛ visual/e2e/a11y تست بنویس. برای این قابلیت، ابتدا acceptance criteria بنویس، سپس migration/schema، service، API، UI، audit، test و documentation را در همان PR انجام بده. & Unit + integration + e2e؛ تست permission، validation، audit و حالت خطا. \\ \hline
277 & empty states & design system را با tokens، RTL، dark/light، responsive و animation کنترل شده بازسازی کن؛ صفحات role-based را polish کن؛ visual/e2e/a11y تست بنویس. برای این قابلیت، ابتدا acceptance criteria بنویس، سپس migration/schema، service، API، UI، audit، test و documentation را در همان PR انجام بده. & Unit + integration + e2e؛ تست permission، validation، audit و حالت خطا. \\ \hline
278 & error states & design system را با tokens، RTL، dark/light، responsive و animation کنترل شده بازسازی کن؛ صفحات role-based را polish کن؛ visual/e2e/a11y تست بنویس. برای این قابلیت، ابتدا acceptance criteria بنویس، سپس migration/schema، service، API، UI، audit، test و documentation را در همان PR انجام بده. & Unit + integration + e2e؛ تست permission، validation، audit و حالت خطا. \\ \hline
279 & toast notification & design system را با tokens، RTL، dark/light، responsive و animation کنترل شده بازسازی کن؛ صفحات role-based را polish کن؛ visual/e2e/a11y تست بنویس. برای این قابلیت، ابتدا acceptance criteria بنویس، سپس migration/schema، service، API، UI، audit، test و documentation را در همان PR انجام بده. & Unit + integration + e2e؛ تست permission، validation، audit و حالت خطا. \\ \hline
280 & command palette & design system را با tokens، RTL، dark/light، responsive و animation کنترل شده بازسازی کن؛ صفحات role-based را polish کن؛ visual/e2e/a11y تست بنویس. برای این قابلیت، ابتدا acceptance criteria بنویس، سپس migration/schema، service، API، UI، audit، test و documentation را در همان PR انجام بده. & Unit + integration + e2e؛ تست permission، validation، audit و حالت خطا. \\ \hline
\end{longtable}
\subsection{P) مدیریت آموزش}

\begin{longtable}{|>{\raggedleft\arraybackslash}p{0.065\textwidth}|>{\raggedleft\arraybackslash}p{0.21\textwidth}|>{\raggedleft\arraybackslash}p{0.42\textwidth}|>{\raggedleft\arraybackslash}p{0.18\textwidth}|}
\hline
\textbf{شماره} & \textbf{قابلیت} & \textbf{مراحل پیاده سازی} & \textbf{تست و پذیرش} \\\hline
\endfirsthead
\hline
\textbf{شماره} & \textbf{قابلیت} & \textbf{مراحل پیاده سازی} & \textbf{تست و پذیرش} \\\hline
\endhead
281 & announcement university-level & announcement/calendar/admin workflows را کامل کن؛ scheduling و expiry را enforce کن؛ email/notification و ICS اضافه کن؛ تست scope، expiry و read receipt بنویس. برای این قابلیت، ابتدا acceptance criteria بنویس، سپس migration/schema، service، API، UI، audit، test و documentation را در همان PR انجام بده. & Unit + integration + e2e؛ تست permission، validation، audit و حالت خطا. \\ \hline
282 & announcement department-level & announcement/calendar/admin workflows را کامل کن؛ scheduling و expiry را enforce کن؛ email/notification و ICS اضافه کن؛ تست scope، expiry و read receipt بنویس. برای این قابلیت، ابتدا acceptance criteria بنویس، سپس migration/schema، service، API، UI، audit، test و documentation را در همان PR انجام بده. & Unit + integration + e2e؛ تست permission، validation، audit و حالت خطا. \\ \hline
283 & announcement course-level & announcement/calendar/admin workflows را کامل کن؛ scheduling و expiry را enforce کن؛ email/notification و ICS اضافه کن؛ تست scope، expiry و read receipt بنویس. برای این قابلیت، ابتدا acceptance criteria بنویس، سپس migration/schema، service، API، UI، audit، test و documentation را در همان PR انجام بده. & Unit + integration + e2e؛ تست permission، validation، audit و حالت خطا. \\ \hline
284 & priority announcement & announcement/calendar/admin workflows را کامل کن؛ scheduling و expiry را enforce کن؛ email/notification و ICS اضافه کن؛ تست scope، expiry و read receipt بنویس. برای این قابلیت، ابتدا acceptance criteria بنویس، سپس migration/schema، service، API، UI، audit، test و documentation را در همان PR انجام بده. & Unit + integration + e2e؛ تست permission، validation، audit و حالت خطا. \\ \hline
285 & scheduled announcement & announcement/calendar/admin workflows را کامل کن؛ scheduling و expiry را enforce کن؛ email/notification و ICS اضافه کن؛ تست scope، expiry و read receipt بنویس. برای این قابلیت، ابتدا acceptance criteria بنویس، سپس migration/schema، service، API، UI، audit، test و documentation را در همان PR انجام بده. & Unit + integration + e2e؛ تست permission، validation، audit و حالت خطا. \\ \hline
286 & expiry announcement & announcement/calendar/admin workflows را کامل کن؛ scheduling و expiry را enforce کن؛ email/notification و ICS اضافه کن؛ تست scope، expiry و read receipt بنویس. برای این قابلیت، ابتدا acceptance criteria بنویس، سپس migration/schema، service، API، UI، audit، test و documentation را در همان PR انجام بده. & Unit + integration + e2e؛ تست permission، validation، audit و حالت خطا. \\ \hline
287 & academic calendar & announcement/calendar/admin workflows را کامل کن؛ scheduling و expiry را enforce کن؛ email/notification و ICS اضافه کن؛ تست scope، expiry و read receipt بنویس. برای این قابلیت، ابتدا acceptance criteria بنویس، سپس migration/schema، service، API، UI، audit، test و documentation را در همان PR انجام بده. & Unit + integration + e2e؛ تست permission، validation، audit و حالت خطا. \\ \hline
288 & exam date & announcement/calendar/admin workflows را کامل کن؛ scheduling و expiry را enforce کن؛ email/notification و ICS اضافه کن؛ تست scope، expiry و read receipt بنویس. برای این قابلیت، ابتدا acceptance criteria بنویس، سپس migration/schema، service، API، UI، audit، test و documentation را در همان PR انجام بده. & Unit + integration + e2e؛ تست permission، validation، audit و حالت خطا. \\ \hline
289 & add/drop deadline & announcement/calendar/admin workflows را کامل کن؛ scheduling و expiry را enforce کن؛ email/notification و ICS اضافه کن؛ تست scope، expiry و read receipt بنویس. برای این قابلیت، ابتدا acceptance criteria بنویس، سپس migration/schema، service، API، UI، audit، test و documentation را در همان PR انجام بده. & Unit + integration + e2e؛ تست permission، validation، audit و حالت خطا. \\ \hline
290 & holiday calendar & announcement/calendar/admin workflows را کامل کن؛ scheduling و expiry را enforce کن؛ email/notification و ICS اضافه کن؛ تست scope، expiry و read receipt بنویس. برای این قابلیت، ابتدا acceptance criteria بنویس، سپس migration/schema، service، API، UI، audit، test و documentation را در همان PR انجام بده. & Unit + integration + e2e؛ تست permission، validation، audit و حالت خطا. \\ \hline
291 & course event & announcement/calendar/admin workflows را کامل کن؛ scheduling و expiry را enforce کن؛ email/notification و ICS اضافه کن؛ تست scope، expiry و read receipt بنویس. برای این قابلیت، ابتدا acceptance criteria بنویس، سپس migration/schema، service، API، UI، audit، test و documentation را در همان PR انجام بده. & Unit + integration + e2e؛ تست permission، validation، audit و حالت خطا. \\ \hline
292 & admin import calendar & announcement/calendar/admin workflows را کامل کن؛ scheduling و expiry را enforce کن؛ email/notification و ICS اضافه کن؛ تست scope، expiry و read receipt بنویس. برای این قابلیت، ابتدا acceptance criteria بنویس، سپس migration/schema، service، API، UI، audit، test و documentation را در همان PR انجام بده. & Unit + integration + e2e؛ تست permission، validation، audit و حالت خطا. \\ \hline
293 & calendar ICS export & announcement/calendar/admin workflows را کامل کن؛ scheduling و expiry را enforce کن؛ email/notification و ICS اضافه کن؛ تست scope، expiry و read receipt بنویس. برای این قابلیت، ابتدا acceptance criteria بنویس، سپس migration/schema، service، API، UI، audit، test و documentation را در همان PR انجام بده. & Unit + integration + e2e؛ تست permission، validation، audit و حالت خطا. \\ \hline
294 & email announcement & announcement/calendar/admin workflows را کامل کن؛ scheduling و expiry را enforce کن؛ email/notification و ICS اضافه کن؛ تست scope، expiry و read receipt بنویس. برای این قابلیت، ابتدا acceptance criteria بنویس، سپس migration/schema، service، API، UI، audit، test و documentation را در همان PR انجام بده. & Unit + integration + e2e؛ تست permission، validation، audit و حالت خطا. \\ \hline
295 & notification announcement & announcement/calendar/admin workflows را کامل کن؛ scheduling و expiry را enforce کن؛ email/notification و ICS اضافه کن؛ تست scope، expiry و read receipt بنویس. برای این قابلیت، ابتدا acceptance criteria بنویس، سپس migration/schema، service، API، UI، audit، test و documentation را در همان PR انجام بده. & Unit + integration + e2e؛ تست permission، validation، audit و حالت خطا. \\ \hline
296 & pinned announcement & announcement/calendar/admin workflows را کامل کن؛ scheduling و expiry را enforce کن؛ email/notification و ICS اضافه کن؛ تست scope، expiry و read receipt بنویس. برای این قابلیت، ابتدا acceptance criteria بنویس، سپس migration/schema، service، API، UI، audit، test و documentation را در همان PR انجام بده. & Unit + integration + e2e؛ تست permission، validation، audit و حالت خطا. \\ \hline
297 & read receipt announcement & announcement/calendar/admin workflows را کامل کن؛ scheduling و expiry را enforce کن؛ email/notification و ICS اضافه کن؛ تست scope، expiry و read receipt بنویس. برای این قابلیت، ابتدا acceptance criteria بنویس، سپس migration/schema، service، API، UI، audit، test و documentation را در همان PR انجام بده. & Unit + integration + e2e؛ تست permission، validation، audit و حالت خطا. \\ \hline
298 & archive announcement & announcement/calendar/admin workflows را کامل کن؛ scheduling و expiry را enforce کن؛ email/notification و ICS اضافه کن؛ تست scope، expiry و read receipt بنویس. برای این قابلیت، ابتدا acceptance criteria بنویس، سپس migration/schema، service، API، UI، audit، test و documentation را در همان PR انجام بده. & Unit + integration + e2e؛ تست permission، validation، audit و حالت خطا. \\ \hline
299 & department admin role & announcement/calendar/admin workflows را کامل کن؛ scheduling و expiry را enforce کن؛ email/notification و ICS اضافه کن؛ تست scope، expiry و read receipt بنویس. برای این قابلیت، ابتدا acceptance criteria بنویس، سپس migration/schema، service، API، UI، audit، test و documentation را در همان PR انجام بده. & Unit + integration + e2e؛ تست permission، validation، audit و حالت خطا. \\ \hline
300 & report to education office & announcement/calendar/admin workflows را کامل کن؛ scheduling و expiry را enforce کن؛ email/notification و ICS اضافه کن؛ تست scope، expiry و read receipt بنویس. برای این قابلیت، ابتدا acceptance criteria بنویس، سپس migration/schema، service، API، UI، audit، test و documentation را در همان PR انجام بده. & Unit + integration + e2e؛ تست permission، validation، audit و حالت خطا. \\ \hline
\end{longtable}
\section{اولویت اجرایی پیشنهادی}
\subsection{فاز ۱ - پایداری و باگ گیری}
\begin{itemize}
\item اجرای pnpm install
\item اجرای pnpm db:generate
\item اجرای pnpm db:migrate
\item اجرای pnpm db:seed
\item اجرای pnpm typecheck
\item اجرای pnpm test
\item اجرای pnpm test:e2e
\item اجرای pnpm build
\item رفع خطاهای build/runtime
\item ثبت همه باگ ها در BUGS.md
\end{itemize}
\subsection{فاز ۲ - اصلاح دیتابیس و امنیت}
\begin{itemize}
\item تست raw SQL constraintها
\item تست partial unique indexها
\item تست anonymous survey/poll
\item تست grade bounds
\item تست permission escalation
\item تست revoked role
\item تست file access
\item تست export audit
\end{itemize}
\subsection{فاز ۳ - تکمیل محصول}
\begin{itemize}
\item Assignment submission
\item Rubric grading
\item Regrade workflow کامل
\item Office hour attendance
\item Task management
\item Timesheet
\item Talent pool
\item Recommendation engine
\end{itemize}
\subsection{فاز ۴ - UI/UX}
\begin{itemize}
\item landing page مدرن
\item dashboardهای حرفه ای
\item animations کنترل شده
\item Jalali calendar
\item responsive tables
\item dark mode کامل
\item accessibility
\end{itemize}
\subsection{۱۰ کار مهم اگر زمان محدود است}
\begin{itemize}
\item pnpm typecheck و رفع همه خطاها
\item pnpm build و رفع runtime/build errors
\item تست کامل migration و seed از صفر
\item تست permission escalation
\item تست anonymous survey/poll
\item تست gradebook و import/export
\item اضافه کردن assignment submission
\item کامل کردن role-based dashboards
\item اضافه کردن UI مدرن تر برای landing/dashboard
\item نوشتن مستندات نهایی و سناریوهای demo
\end{itemize}

\section{پرامپت جامع آماده برای Claude}
متن این بخش را می توانی مستقیم به Claude بدهی. هدف این پرامپت این است که Claude پروژه را هم از نظر باگ، هم از نظر معماری، هم از نظر تست و هم از نظر UI/UX بررسی و اصلاح کند.

\begin{framed}
\textbf{نقش تو:} تو یک \lr{Senior Full-Stack Engineer}، \lr{Security Reviewer}، \lr{QA Lead} و \lr{Product Architect} هستی. مخزن زیر را کامل بررسی کن:
\begin{latin}
\url{https://github.com/Ramin-Buzarpur/TAFlow}
\end{latin}

\textbf{هدف:} پروژه \lr{TAFlow} باید از یک MVP سنگین به یک سامانه دانشگاهی پایدار، تست شده، امن، role-based، RTL و قابل ارائه تبدیل شود.

\textbf{قواعد کار:}
\begin{enumerate}
\item اول پروژه را locally clone کن، dependencies را فقط با pnpm نصب کن، نسخه ها را مطابق packageManager و lockfile نگه دار.
\item قبل از هر تغییر، این دستورها را اجرا کن و خطاها را در \lr{BUGS.md} ثبت کن:
\begin{latin}
\begin{verbatim}
pnpm install
pnpm db:generate
docker compose up -d
pnpm db:migrate
pnpm db:seed
pnpm typecheck
pnpm test
pnpm test:e2e
pnpm build
\end{verbatim}
\end{latin}
\item هیچ route حساسی نباید صرفا به client state اعتماد کند. همه permissionها باید server-side باشند.
\item همه تغییرات حساس باید audit شوند: role changes، application status، grade changes، file downloads، exports، certificate actions و security events.
\item برای هر باگ یا feature، migration/schema، service، API، UI، validation، permission، audit، unit test، integration/e2e test و documentation را در یک مسیر منظم انجام بده.
\item اول اصلاحات حیاتی را انجام بده: package/version stability، Prisma/migration integrity، RBAC، permission escalation tests، survey/poll anonymity، gradebook security، file access security و seed idempotency.
\item سپس قابلیت های محصولی را اضافه کن: assignment submission، rubric grading، regrade workflow، office hour attendance، task management، timesheet، talent pool، recommendation engine، certificate verification/revocation، notification channels و reports.
\item UI/UX را بعد از پایدار شدن منطق core مدرن کن: landing page، role-based dashboards، dark/light mode، RTL کامل، responsive tables، skeleton/empty/error states، animations کنترل شده و accessibility.
\item برای UI مدرن، از الگوهای Magic UI، Aceternity UI، Motion، AOS و scroll-based interactions الهام بگیر، اما performance و accessibility را خراب نکن. هر animation باید قابل خاموش شدن با \lr{prefers-reduced-motion} باشد.
\item خروجی نهایی باید شامل PR summary، فهرست فایل های تغییرکرده، migrationها، تست های اضافه شده، دستورات اجرا، مشکلات باقی مانده و demo checklist باشد.
\end{enumerate}

\textbf{الگوی اجرای هر قابلیت:}
\begin{enumerate}
\item \lr{Acceptance Criteria}
\item \lr{Data Model / Migration}
\item \lr{Service Layer}
\item \lr{API Route}
\item \lr{Permission / RBAC}
\item \lr{Audit / SecurityEvent}
\item \lr{UI/UX}
\item \lr{Unit Test}
\item \lr{Integration/E2E Test}
\item \lr{Documentation}
\end{enumerate}

\textbf{شرط پایان کار:} قبل از پایان، این دستورها باید بدون خطا pass شوند:
\begin{latin}
\begin{verbatim}
pnpm typecheck
pnpm test
pnpm test:e2e
pnpm build
\end{verbatim}
\end{latin}
\end{framed}

\section{معیار پذیرش نهایی پروژه}
\begin{itemize}
\item نصب از صفر روی ویندوز و لینوکس باید با دستورهای README قابل تکرار باشد.
\item migrate و seed از دیتابیس خالی باید بدون خطا انجام شود.
\item typecheck، test، test:e2e و build باید بدون خطا pass شوند.
\item تمام نقش ها شامل Student، TA، Head TA، Professor و Education Admin باید dashboard و flow مشخص داشته باشند.
\item هیچ route حساس نباید بدون permission server-side عمل کند.
\item گزارش \lr{BUGS.md}، \lr{SECURITY\_NOTES.md}، \lr{QA\_CHECKLIST.md} و \lr{DEMO\_SCRIPT.md} باید کامل شود.
\item UI باید RTL، responsive، دارای loading/empty/error state و سازگار با dark/light mode باشد.
\end{itemize}
\end{document}
